\section{Status, Records, Coherence}\label{sec:status-records-coherence}

This chapter treats records as part of the mathematical object rather
than as commentary about it.  A formed fact has a claim record, a
status, and a place in a coherence ledger that records how it
persists, transports, composes and supports explanation; the laws of
the chapter govern that bookkeeping.  In the running example of
\Cref{sec:intro}, the record of the split pair \((1,2)\), with its
status and the repair chosen, is part of the object these laws
describe.

The fourteen laws run from shared objects to counterfactual use.
Objectivity as Common Quotient identifies when several presentations
share an object; Lawhood Descent tests when a law factors through the
available quotient; Object Persistence tests when an object survives
transport; and Fact / Record Formation says what it means for a fact
or record to form.  Reflexive Nonclosure, Interface Mediation,
Probability as Stable Unresolved Fiber and Causality as Stable
Intervention Transport describe how records behave under self-audit,
interfaces, unresolved fibers and interventions.  The remaining laws
concern coherence: Moduli / Contingency separates stable parameters
from contingent choices, Conservation as Orbit Descent reads
conservation as a quotient law on orbits, Composability and
Presentation Invariance say when records survive composition and change
of presentation, Explanation as Compression with Descent treats
explanation as compression that still factors through the declared
quotient, and Counterfactual Legitimacy says when a counterfactual
claim is lawful rather than an overread.

The chapter uses only the vocabulary of \Cref{sec:apparatus}.  The
summary table at the end groups the laws into three reading groups,
which do not follow document order: status (\Fref{F17}, \Fref{F18},
\Fref{F22}, \Fref{F23}), records (\Fref{F19}, \Fref{F20}, \Fref{F21},
\Fref{F25}) and coherence (\Fref{F26}--\Fref{F31}).  The chapter's axes
are status, records and interfaces (\Cref{subsec:reading-guide});
coherence is a reading group, not an axis.

\subsection{Objectivity as Common Quotient [\texorpdfstring{\Fref{F17}}{F17}]}\label{sec:status-records-coherence:objectivity}

\paragraph{Reading the statement.} The symbols are those of \Cref{thm:F17} below.
Here ``\(r\) descends through \(q\)'' means that \(r=\bar r\circ q\) for
some map \(\bar r\), as in \Cref{sec:apparatus}.
\(\operatorname{CommonAcrossAccesses}(q,r)\) says that \(r\) descends
through every access quotient \(q_i\), and
\(\operatorname{PublicDeterminesObject}(p,r)\) says that \(r\) descends
through the public quotient \(p\).
\(\operatorname{AccessObjectivityDefectEmpty}(q_i,r)\) says that no two
histories with the same \(q_i\)-class have different \(r\)-values.
\(\operatorname{PublicToObjectDefect}(p,r)\) is the set of pairs with
the same \(p\)-class and different \(r\)-values, and
\(\operatorname{ObstructionEmpty}\) says that this set is empty.
\(\operatorname{CommonQuotientVisible}(q,p)\) says that \(p\) itself
descends through every \(q_i\). The objective extension
\(\operatorname{objectiveExtension}(p,r)\) is the pairing
\(h\mapsto(p(h),r(h))\); it \emph{retains} \(p\), and \emph{carries}
\(r\), when \(p\), respectively \(r\), descends through it.
The status partition
\(\operatorname{ExactlyOneObjectivityStatus}(p,r)\) says that exactly
one of four statuses holds, according to whether \(r\) descends through
\(p\) and whether \(p\) descends through \(r\): exact common quotient when
both hold, coarse objective object when only \(r\) descends through \(p\),
private overrefinement when only \(p\) descends through \(r\), and
incomparable object claim when neither does.

\begin{theorem}[Objectivity as Common Quotient]\label{thm:F17}
Let \(H\) be a carrier with a family of access quotients
\(q_i:H\to Q_i\) indexed by \(i\in I\), a public quotient
\(p:H\to P\), and a target readout \(r:H\to R\).  Assume each
\(q_i\) is surjective, \(p\) is surjective, and the family satisfies
the common-quotient universal property
\(\operatorname{CommonQuotientUniversal}(q,p)\): \(p\) factors
through every \(q_i\), and any map \(s:H\to S\) that factors through every
\(q_i\) also factors through \(p\).  Then
\begin{align*}
  &\operatorname{CommonAcrossAccesses}(q,r)
    \iff \operatorname{PublicDeterminesObject}(p,r),\\
  &\operatorname{CommonAcrossAccesses}(q,r)\\
  &\qquad{}\iff
    \forall i\in I,\,
      \operatorname{AccessObjectivityDefectEmpty}(q_i,r),\\
  &\operatorname{PublicDeterminesObject}(p,r)\\
  &\qquad{}\iff
    \operatorname{ObstructionEmpty}
      \bigl(\operatorname{PublicToObjectDefect}(p,r)\bigr),
\end{align*}
together with
\begin{align*}
  &\operatorname{CommonQuotientVisible}(q,p),\\
  &\operatorname{RetainsPublicQuotient}
      \bigl(p,\operatorname{objectiveExtension}(p,r)\bigr),\\
  &\operatorname{CarriesObject}
      \bigl(\operatorname{objectiveExtension}(p,r),r\bigr),
\end{align*}
and the status partition
\[
  \operatorname{ExactlyOneObjectivityStatus}(p,r).
\]
In ordinary terms: \(p\) factors through every \(q_i\); the pairing
\(\langle p,r\rangle\) satisfies \(p=\mathrm{pr}_1\circ\langle p,r\rangle\) and
\(r=\mathrm{pr}_2\circ\langle p,r\rangle\); and exactly one of the following holds:
both \(r\) factors through \(p\) and \(p\) through \(r\) (exact common quotient);
only the first (coarse objective object); only the second (private
overrefinement); neither (incomparable).
\end{theorem}
\begin{proof}
Assume the surjectivity of each \(q_i\), the surjectivity of \(p\),
and \(\operatorname{CommonQuotientUniversal}(q,p)\).  We prove the
seven conjuncts in turn.

\emph{(1) Common across accesses iff public determines object.}
The universal property supplies comparison maps
\[
  p=\pi_i\circ q_i\quad(i\in I),
  \qquad
  s=\lambda\circ p
  \quad\text{whenever } s=\mu_i\circ q_i \text{ for every } i\in I .
\]
Thus \(p\) is the finest quotient that factors through every \(q_i\).
If \(r\) is common across the access quotients, then \(r\) factors through
every \(q_i\) by definition, hence through
\(p\), by the universal property applied to the map \(r\).  Conversely, if \(r\) factors through \(p\), it factors through
every \(q_i\), since \(p=\pi_i\circ q_i\).  So a readout \(r\) is common
across the access quotients exactly when it is constant on the common
public fibers:
\[
  \forall h,h'\in H,\quad
  p(h)=p(h')\Longrightarrow r(h)=r(h').
\]
Since \(p\) is surjective, this is equivalent to the existence of
\(\bar r:P\to R\) with
\[
  r=\bar r\circ p,
\]
which is precisely \(\operatorname{PublicDeterminesObject}(p,r)\).

\emph{(2) Common across accesses iff every access defect is empty.}
For \(i\in I\), the access objectivity defect is the set
\[
  \Delta_i(q_i,r)
  =
  \{(h,h')\in H\times H:
      q_i(h)=q_i(h')\ \text{and}\ r(h)\neq r(h')\}.
\]
Since \(q_i\) is surjective, the descent criterion of
\Cref{sec:apparatus} shows that \(\Delta_i(q_i,r)=\emptyset\) is
equivalent to \(r=\bar r_i\circ q_i\) for some \(\bar r_i\); requiring
this for every \(i\in I\) is \(\operatorname{CommonAcrossAccesses}(q,r)\),
giving the second equivalence.

\emph{(3) Public determines object iff the public-to-object
obstruction is empty.}
The public-to-object defect is
\[
  \Delta_P(p,r)
  =
  \{(h,h')\in H\times H:
      p(h)=p(h')\ \text{and}\ r(h)\neq r(h')\}.
\]
By the descent criterion, since \(p\) is surjective, its emptiness is
equivalent to \(r=\bar r\circ p\) for some \(\bar r:P\to R\), that is, to
\(\operatorname{PublicDeterminesObject}(p,r)\).

\emph{(4) Common quotient visible.}
The same universal property exhibits \(p\) as a declared common
quotient of the family \((q_i)_{i\in I}\):
\[
  \operatorname{CommonQuotientUniversal}(q,p)
  \Longrightarrow
  \operatorname{CommonQuotientVisible}(q,p).
\]
Here visibility is the first half of the universal property: \(p\)
descends through every \(q_i\), \(p=\pi_i\circ q_i\).

\emph{(5) The objective extension retains the public quotient.}
Let
\[
  \operatorname{objectiveExtension}(p,r)
  :=
  \langle p,r\rangle:H\longrightarrow P\times R,
  \qquad
  h\longmapsto (p(h),r(h)).
\]
Projection to the first coordinate recovers the public quotient:
\[
  \operatorname{pr}_P\circ\langle p,r\rangle=p.
\]
Therefore the objective extension retains \(p\).

\emph{(6) The objective extension carries the object.}
Projection to the second coordinate recovers the object readout:
\[
  \operatorname{pr}_R\circ\langle p,r\rangle=r.
\]
Thus the same extension carries the object \(r\).

\emph{(7) Exactly one objectivity status.}
The four status predicates compare the object \(r\) with the public
quotient \(p\) in both directions: whether \(r\) descends through \(p\)
(the public quotient determines the object) and whether \(p\) descends
through \(r\) (the object determines the public quotient).  The cases
are
\[
\begin{array}{c|c|c}
\text{status} & r\ \text{descends through}\ p
  & p\ \text{descends through}\ r\\
\hline
\text{exact common quotient} & \text{yes} & \text{yes}\\
\text{coarse objective object} & \text{yes} & \text{no}\\
\text{private overrefinement} & \text{no} & \text{yes}\\
\text{incomparable object claim} & \text{no} & \text{no}
\end{array}
\]
Every configuration of the two Boolean tests falls into one row, and
no configuration falls into two rows.  Hence
\(\operatorname{ExactlyOneObjectivityStatus}(p,r)\).
\end{proof}

\paragraph{Nonclaims.} The theorem does not determine which presentations are declarable and
does not enumerate the kinds of objects that may be shared.  The
closest physical analogy is quantum Darwinism, in which
objectivity emerges as redundant agreement of the environment's
records across many access channels \citep{Zurek2009QuantumDarwinism};
the layer-agnostic statement above is independent of that
substrate.

\subsection*{Layer instantiations}\label{layer-inst:F17}

\begingroup\small
\begin{description}
\item[\emph{Fixed-effect meta-analysis (statistics).}] \textit{Analogy.}
Under a declared fixed-common-effect model, a fixed-effect
meta-analysis combines independent study-level effect-size
estimates by inverse-variance weighting to produce a pooled
consensus estimate computed from the study-level estimates; when the
model holds, every study estimates the same common effect, and the
pooled estimate is its public estimator.  Cochran's \(Q\) (and the later \(I^2\) index of Higgins and Thompson) measures the
\Fref{F17} access-objectivity defect: heterogeneity is the witness that
the declared common-effect model is violated
\citep{HedgesOlkin1985MetaAnalysis}.

\end{description}
\endgroup


\subsection{Lawhood Descent [\texorpdfstring{\Fref{F18}}{F18}]}\label{sec:status-records-coherence:lawhood-descent}

\paragraph{Reading the statement.}
A candidate law here is a predicate or readout \(P:H\to V\). It is
\emph{lawful} at the layer when it is invariant under the symmetries of the
access quotient: \(P\circ\sigma=P\) for every bijection \(\sigma:H\to H\)
with \(\accessQ\circ\sigma=\accessQ\). For a relation or a transition rule,
lawfulness means factoring through \(\accessQ\) in each argument; invariance
under such \(\sigma\) is weaker there, since every such \(\sigma\) preserves
the equality relation on \(H\), which need not factor.

\begin{theorem}[Lawhood Descent]\label{thm:F18}
Let \(\accessQ:H\to Q\) be a surjective access quotient (the available access of
a formed closure in the intended reading).  A declared candidate law \(P:H\to V\) satisfies \(P\circ\sigma=P\) for
every bijection \(\sigma:H\to H\) with \(\accessQ\circ\sigma=\accessQ\) (that is,
\(P\) is lawful at the layer) if and only if \(P\) factors through \(\accessQ\),
that is, \(P(h)=P(h')\) whenever \(\accessQ(h)=\accessQ(h')\).
\end{theorem}
\begin{proof}
If \(P=\bar P\circ\accessQ\) and \(\accessQ\circ\sigma=\accessQ\), then
\(P\circ\sigma=\bar P\circ\accessQ\circ\sigma=\bar P\circ\accessQ=P\).
Conversely, let \(P\) be lawful and \(\accessQ(h)=\accessQ(h')\) with
\(h\ne h'\). The transposition \(\sigma\) exchanging \(h\) and \(h'\) is a
bijection with \(\accessQ\circ\sigma=\accessQ\), so
\(P(h')=P(\sigma(h))=P(h)\). Hence \(P\) is constant on access classes;
define \(\bar P(\accessQ(h))=P(h)\), which surjectivity defines on all of
\(Q\).  Constancy makes this well-defined, and
\(P=\bar P\circ\accessQ\).  Thus lawhood is precisely quotient descent.
\end{proof}

\paragraph{Nonclaims.} The theorem does not enumerate lawful candidates and does not
construct the access quotient; that is the role of Quotientality.

\subsection*{Layer instantiations}\label{layer-inst:F18}

\begingroup\small
\begin{description}
\item[\emph{Wilson loops in gauge theory (physics).}] \textit{Special case.}
In Yang--Mills theory, a candidate physical observable is a
lawful observable at the gauge-quotient iff it is
gauge-invariant.  Wilson loops --- traces of the holonomy around
closed loops --- are the canonical gauge-invariant
observables, descending through the gauge-orbit projection
\(\mathcal A\to\mathcal A/\mathcal G\)
\citep{Wilson1974ConfinementQuarks}.

\item[\emph{Parnas information hiding (software engineering).}] \textit{Analogy.}
Parnas's principle of information hiding declares that a module
exposes a public interface and hides its implementation behind
that interface.  A client predicate is a lawful use of the
module iff it depends only on the declared public operations;
predicates that read hidden state are overreads.  The interface
is the \Fref{F18} access quotient through which lawful client claims
descend
\citep{Parnas1972InformationHiding}.

\item[\emph{Bisimulation invariance in process calculi
(concurrency theory).}] \textit{Special case.}
In Milner--Park bisimulation theory, every Hennessy--Milner
modal-logic formula is bisimulation-invariant, and on image-finite
labelled transition systems two processes are bisimilar iff they
satisfy the same formulas.  A process property is lawful iff it is
bisimulation-invariant, and the bisimulation quotient on the labelled
transition system is the \Fref{F18} access quotient through which
lawful process-level laws, the Hennessy--Milner formulas among them,
descend
\citep{Milner1989CommunicationConcurrency}.

\item[\emph{Congruences and quotient algebras (universal
algebra).}] \textit{Special case.}
A predicate on an algebra \(A\) descends to a well-defined predicate
on the quotient \(A/{\sim}\) iff it is \(\sim\)-invariant, and the
operations of \(A\) descend to \(A/{\sim}\) iff \(\sim\) is a
congruence (compatible with all operations); for a congruence
\(\sim\), laws over the quotient algebra correspond exactly to
\(\sim\)-invariant predicates on the original algebra
\citep{Gratzer1979UniversalAlgebra}.

\end{description}
\endgroup


\subsection{Object Persistence [\texorpdfstring{\Fref{F19}}{F19}]}\label{sec:status-records-coherence:object-persistence}

\paragraph{Reading the statement.} The symbols are those of \Cref{thm:F19} below.
\(\operatorname{Persists}(q_i,\gamma,q_j)\) says that the later class
\(q_j(\gamma h)\) depends only on the earlier class \(q_i(h)\), that
is, \(q_j\circ\gamma\) descends through \(q_i\). The persistence
obstruction, also called the split obstruction, is the set of pairs
with the same \(q_i\)-class whose transports have different
\(q_j\)-classes; the \emph{ObstructionEmpty} predicates say that the
named set is empty. \(\operatorname{ViablePersistence}\) says that every
transported class is viable, and the dissolution obstruction is the set
of histories whose transported class is not viable.
\(\operatorname{RelationValuedPersistence}(q_i,q_j,C)\) says that
whether \(C(h,k)\) holds depends only on the classes \(q_i(h)\) and
\(q_j(k)\); its obstruction is the set of pairs of pairs with the same
classes on which \(C\) differs.
\(\operatorname{IdentityHolonomyCoherent}\) says that
\(\operatorname{compBar}=\bar\eta\circ\bar\gamma\), and the holonomy
defect is the set of classes where the two differ. The comparison
status compares the split obstruction with the merge obstruction, the
pairs merged by the transport that \(q_i\) separates.  The persistence
extension is the pairing \(h\mapsto(q_i(h),q_j(\gamma h))\); it
\emph{retains} the old identity \(q_i\), and \emph{carries} the later
identity \(q_j\circ\gamma\), when that map descends through it.

\begin{theorem}[Object Persistence]\label{thm:F19}
Let \(H_i,H_j\) be carriers with access quotients
\(q_i:H_i\to Q_i\) and \(q_j:H_j\to Q_j\), and let
\(\gamma:H_i\to H_j\) be a declared transport.  Assume \(q_i\) is surjective.  Let
\(\operatorname{viable}:Q_j\to\mathrm{Prop}\) be a viability
predicate, \(C:H_i\times H_j\to\mathrm{Prop}\) a relation predicate,
let \(Q_k\) be a declared target set, and let
\[
  \bar\gamma:Q_i\to Q_j,\qquad
  \bar\eta:Q_j\to Q_k,\qquad
  \operatorname{compBar}:Q_i\to Q_k
\]
be arbitrary holonomy data (no compatibility between
\(\operatorname{compBar}\) and \(\bar\eta\circ\bar\gamma\) is
assumed).  Then
\begin{align*}
  &\operatorname{Persists}(q_i,\gamma,q_j)
    \iff
    \operatorname{PersistenceObstructionEmpty}(q_i,\gamma,q_j)\\
  &\qquad{}\iff
    \operatorname{SplitObstructionEmpty}(q_i,\gamma,q_j),\\
  &\operatorname{ViablePersistence}(\gamma,q_j,\operatorname{viable})
    \iff
    \operatorname{DissolutionObstructionEmpty}
      (\gamma,q_j,\operatorname{viable}),\\
  &\operatorname{RelationValuedPersistence}(q_i,q_j,C)
    \iff
    \operatorname{RelationPersistenceObstructionEmpty}(q_i,q_j,C),\\
  &\operatorname{IdentityHolonomyCoherent}
      (\bar\gamma,\bar\eta,\operatorname{compBar})\\
  &\qquad{}\iff
    \operatorname{IdentityHolonomyDefectEmpty}
      (\bar\gamma,\bar\eta,\operatorname{compBar}),
\end{align*}
where \(\operatorname{ViablePersistence}:\Longleftrightarrow\forall h,\
\operatorname{viable}(q_j(\gamma h))\);
\(\operatorname{RelationValuedPersistence}:\Longleftrightarrow C=\bar C\circ(q_i\times
q_j)\) for some \(\bar C\); and
\(\operatorname{IdentityHolonomyCoherent}:\Longleftrightarrow\operatorname{compBar}=\bar\eta\circ\bar\gamma\);
together with the comparison-status partition
\[
  \operatorname{ExactlyOnePersistenceComparisonStatus}
    (q_i,\gamma,q_j),
\]
and the retention and carrying conditions
\begin{align*}
  &\operatorname{RetainsOldIdentity}
      \bigl(q_i,\operatorname{persistenceExtension}(q_i,\gamma,q_j)\bigr),\\
  &\operatorname{CarriesLaterIdentity}
      \bigl(\operatorname{persistenceExtension}(q_i,\gamma,q_j),
        \gamma,q_j\bigr).
\end{align*}
If moreover \(\eta:H_j\to H_k\) and \(q_k:H_k\to Q_k\) satisfy
\(\bar\gamma\circ q_i=q_j\circ\gamma\), \(\bar\eta\circ q_j=q_k\circ\eta\) and
\(\operatorname{compBar}\circ q_i=q_k\circ\eta\circ\gamma\), then
\(\operatorname{IdentityHolonomyCoherent}(\bar\gamma,\bar\eta,\operatorname{compBar})\)
holds; hence, when the holonomy defect is nonempty, one of the three
displayed commutation equations fails.
Here the persistence and split obstructions are the same set
\(\Delta_{\mathrm{split}}=\{(h,h'):q_i(h)=q_i(h'),\ q_j(\gamma h)\ne q_j(\gamma h')\}\), and
\(\operatorname{Persists}(q_i,\gamma,q_j)\) says that \(q_j\circ\gamma\)
factors through \(q_i\).
Here \(\operatorname{ExactlyOnePersistenceComparisonStatus}\) says that exactly
one of exact endurance (\(\Delta_{\mathrm{split}}=\Delta_{\mathrm{merge}}=\emptyset\)),
persistent with merge (only \(\Delta_{\mathrm{merge}}\ne\emptyset\)), identity
branching (only \(\Delta_{\mathrm{split}}\ne\emptyset\)) and identity rewrite (both
nonempty) holds, where \(\Delta_{\mathrm{merge}}=\{(h,h'):q_j(\gamma h)=q_j(\gamma
h'),\ q_i(h)\ne q_i(h')\}\).
\end{theorem}
\begin{proof}
Assume the surjectivity of \(q_i\), the predicates
\(\operatorname{viable}\) and \(C\), and the holonomy data
\(\bar\gamma,\bar\eta,\operatorname{compBar}\).  We prove the eight
conjuncts in turn. Part (3) proves the comparison-status partition, which the statement lists after the holonomy line, and part (6) also proves the added holonomy clause.

\emph{(1) Persists iff the persistence obstruction is empty.}
The persistence obstruction is
\[
  \Delta_{\mathrm{pers}}(q_i,\gamma,q_j)
  =
  \{(h,h')\in H_i\times H_i:
      q_i(h)=q_i(h')\ \text{and}\
      q_j(\gamma h)\ne q_j(\gamma h')\}.
\]
Since \(q_i\) is surjective, the descent criterion of
\Cref{sec:apparatus} applied to \(q_j\circ\gamma\) shows that its
emptiness is equivalent to the existence of a map
\(\gamma^{\flat}:Q_i\to Q_j\) satisfying
\[
  \gamma^{\flat}\circ q_i=q_j\circ\gamma,
\]
which is \(\operatorname{Persists}(q_i,\gamma,q_j)\).

\emph{(2) Split-obstruction empty iff persists.}
The split obstruction is the same pair set written in the comparison
vocabulary:
\[
  \Delta_{\mathrm{split}}(q_i,\gamma,q_j)
  =
  \Delta_{\mathrm{pers}}(q_i,\gamma,q_j).
\]
Thus \(\Delta_{\mathrm{split}}=\emptyset\) iff the persistence
obstruction is empty, and by (1) this is equivalent to persistence.

\emph{(3) Exactly-one persistence comparison status.}
The comparison statuses are decided by two split-pair defects on the
declared continuation \(\gamma\): the persistence/split obstruction
\(\Delta_{\mathrm{split}}\) of case~(1), and the merge obstruction
\[
  \Delta_{\mathrm{merge}}(q_i,\gamma,q_j)
  =
  \{(h,h')\in H_i\times H_i:
      q_j(\gamma h)=q_j(\gamma h')\ \text{and}\ q_i(h)\ne q_i(h')\}.
\]
The four statuses partition the configurations by the joint
\(\Delta_{\mathrm{split}}\)\textendash{}\(\Delta_{\mathrm{merge}}\)
emptiness pattern:
\[
\begin{array}{c|c|c}
\text{status} & \Delta_{\mathrm{split}}=\emptyset
  & \Delta_{\mathrm{merge}}=\emptyset\\
\hline
\text{exact endurance} & \text{yes} & \text{yes}\\
\text{persistent with merge} & \text{yes} & \text{no}\\
\text{identity branching} & \text{no} & \text{yes}\\
\text{identity rewrite} & \text{no} & \text{no}
\end{array}
\]
Every configuration of the two Boolean defect tests falls into
exactly one row.  Existence is by case enumeration; exclusivity is by
definitional disjointness of the four status predicates.  The
dissolution and relation-valued obstructions of cases~(4) and~(5) are
separate conditions and do not enter this partition.

\emph{(4) Viable persistence iff the dissolution obstruction is
empty.}
The dissolution obstruction is
\[
  \Delta_{\mathrm{diss}}(\gamma,q_j,\operatorname{viable})
  =
  \{h\in H_i:
      \neg\,\operatorname{viable}(q_j(\gamma h))\}.
\]
Viable persistence says that every transported \(q_j\)-class lies in
the viable region:
\[
  \forall h\in H_i,\quad
  \operatorname{viable}(q_j(\gamma h)).
\]
This is precisely \(\Delta_{\mathrm{diss}}=\emptyset\).

\emph{(5) Relation-valued persistence iff its obstruction is empty.}
The relation-persistence obstruction is
\[
  \Delta_C(q_i,q_j,C)
  =
  \left\{
  \begin{array}{l}
    ((h,k),(h',k'))\in(H_i\times H_j)^2:\\
    q_i(h)=q_i(h'),\ q_j(k)=q_j(k'),\
    C(h,k)\not\iff C(h',k')
  \end{array}
  \right\}.
\]
Emptiness of \(\Delta_C\) says that the truth of \(C(h,k)\) depends only
on the pair of classes \((q_i(h),q_j(k))\).  Setting \(\bar C:=\bot\) off the image of \(q_i\times q_j\), this holds exactly when \(C\) descends to a
well-defined relation on \(Q_i\times Q_j\), which is relation-valued
persistence; this part uses no surjectivity, and surjectivity of \(q_j\) is not assumed.

\emph{(6) Identity holonomy coherent iff the holonomy defect is
empty.}
The holonomy defect is
\[
  \Delta_{\mathrm{hol}}(\bar\gamma,\bar\eta,\operatorname{compBar})
  =
  \{a\in Q_i:
      \operatorname{compBar}(a)\ne\bar\eta(\bar\gamma(a))\}.
\]
Thus \(\Delta_{\mathrm{hol}}=\emptyset\) iff
\[
  \forall a\in Q_i,\quad
  \operatorname{compBar}(a)=\bar\eta(\bar\gamma(a)),
\]
which is the pointwise form of
\(\operatorname{compBar}=\bar\eta\circ\bar\gamma\), i.e. identity
holonomy coherence. Under the descent equations of the added clause,
\(\operatorname{compBar}(q_ih)=q_k\eta\gamma h=\bar\eta(q_j\gamma h)
=\bar\eta\bar\gamma(q_ih)\) for every \(h\), and \(q_i\) is surjective, so
the defect is empty.

\emph{(7) The persistence extension retains the old identity.}
Define the persistence extension as the joint quotient
\[
  \operatorname{persistenceExtension}(q_i,\gamma,q_j)
  :=
  \langle q_i,q_j\circ\gamma\rangle:H_i\to Q_i\times Q_j.
\]
Projection to the first coordinate recovers the old identity:
\[
  \operatorname{pr}_{Q_i}\circ
  \langle q_i,q_j\circ\gamma\rangle=q_i.
\]
Therefore the persistence extension retains \(q_i\).

\emph{(8) The persistence extension carries the later identity.}
Projection to the second coordinate recovers the transported later
identity:
\[
  \operatorname{pr}_{Q_j}\circ
  \langle q_i,q_j\circ\gamma\rangle=q_j\circ\gamma.
\]
Thus the same extension carries \(q_j\) along the declared transport
\(\gamma\).
\end{proof}

\paragraph{Nonclaims.} The theorem does not classify persistent objects or say which
transports should be declared.

\subsection*{Layer instantiations}\label{layer-inst:F19}

\begingroup\small
\begin{description}
\item[\emph{Mendelian inheritance (genetics).}] \textit{Analogy.}
Through meiotic segregation and fertilisation, parental allele
identity at a locus survives into the offspring's diploid
genotype.  The \Fref{F19} reading: allele identity persists when no
mutation or non-disjunction breaks the descent; mutation and
non-disjunction events are the persistence obstruction
\citep{Mendel1866VersuchePflanzenHybriden}.

\item[\emph{Pure-state identity under unitary evolution (quantum
mechanics).}] \textit{Special case.}
For a closed quantum system, time evolution is the unitary
\(U(t)=\exp(-iHt/\hbar)\); the projective ray
\([|\psi\rangle]\) (state modulo global phase) evolves through the
descended projective evolution.  Pure-state identity persists:
\(U(t)(e^{i\phi}\psi)=e^{i\phi}U(t)\psi\), so the persistence
obstruction is empty.  For decoherence, take unit state vectors, fix an initial environment state
and a unitary evolution of the combined system. With the earlier quotient
the system's projective ray and the later quotient its reduced density
operator, corestricted to the reduced states it attains, global phase cancels, so the persistence obstruction is empty.
When the reduced state is mixed, it fails
\(\operatorname{viable}(\rho):\Leftrightarrow\rho^2=\rho\) and contributes to
the dissolution obstruction \citep{vonNeumann1932MathematischeGrundlagen}.

\end{description}
\endgroup


\subsection{Fact / Record Formation [\texorpdfstring{\Fref{F20}}{F20}]}\label{sec:status-records-coherence:fact-record}

\begin{definition}[Record formation]\label{def:record-formation}
Let \(H_0,H_1\) be history sets (in the intended reading, carriers of a
formed closure \(\closureForm(A)\)), with continuation
\(\gamma:H_0\to H_1\) and record readout \(\rho_1:H_1\to R_1\). The
\emph{record} of a history \(h\in H_0\) is its record signature
\(s(h)=\rho_1(\gamma(h))\), and an event readout \(\varepsilon:H_0\to E\) is a
\emph{recorded fact} when \(\varepsilon=\bar\varepsilon\circ s\) for some
\(\bar\varepsilon:R_1\to E\).
\end{definition}

\paragraph{Reading the statement.} The theorem makes
\Cref{def:record-formation} precise for readouts: the record is the
signature \(\rho_1\circ\gamma\).
\(\operatorname{EventRecorded}(\varepsilon,\gamma,\rho_1)\) says that
the event \(\varepsilon\) descends through the record signature
\(\rho_1\circ\gamma\), so the record determines the event.
\(\operatorname{RecordStable}(\rho_s,\theta,\rho_t)\) says that
\(\rho_t\circ\theta\) descends through \(\rho_s\), so the transport
respects records. \(\operatorname{FactValueStable}\) says that
\(d_t\circ\bar\theta=u\circ d_s\), so the fact value moves by the
declared update. \(\operatorname{StableFact}\) is the conjunction of these
three conditions.  They describe the transport of one recorded fact when
\(H_s=H_1\), \(\rho_s=\rho_1\), \(E_s=E\) and \(d_s=\bar\varepsilon\), the
descended event readout of \Cref{def:record-formation}; the last clause
of the theorem treats this case.  \(\operatorname{EventRecorded}(\varepsilon,\gamma,\rho_1)\)
holds exactly when \(\varepsilon\) is a recorded fact in the sense of
\Cref{def:record-formation}.
\(\operatorname{RetainsRecord}\) and \(\operatorname{CarriesEvent}\) say
that the record signature, respectively the event, descends through
the record-event extension. \(\operatorname{IdempotentPackaging}(C)\)
says that packaging an already packaged history changes nothing. The
packaging clause is independent of the other data: it is the standard fact
that an idempotent map is a retraction onto its fixed set, included because a
packaged record must be unchanged by repackaging; \(H\) denotes the auxiliary
carrier only in this clause.

\begin{theorem}[Fact / Record Formation]\label{thm:F20}
Let \(H_0,H_1,H_s,H_t\) be history spaces with event readout
\(\varepsilon:H_0\to E\), continuation \(\gamma:H_0\to H_1\),
record readouts
\[
  \rho_1:H_1\to R_1,\qquad
  \rho_s:H_s\to R_s,\qquad
  \rho_t:H_t\to R_t,
\]
transport \(\theta:H_s\to H_t\), a declared record-level transport
\(\bar\theta:R_s\to R_t\), fact-value maps
\(d_s:R_s\to E_s\), \(d_t:R_t\to E_t\), update
\(u:E_s\to E_t\), and a packaging map \(C:H\to H\) on an auxiliary
carrier \(H\).  Define the record signature
\[
  s:=\rho_1\circ\gamma:H_0\to R_1,
\]
and assume that \(s\) and \(\rho_s\) are surjective.  Define
\begin{align*}
  \mathcal O_{\mathrm{ev}}
  &:=
  \{(h,h')\in H_0^2:
      \rho_1(\gamma h)=\rho_1(\gamma h')
      \text{ and } \varepsilon(h)\ne\varepsilon(h')\},\\
  \mathcal O_{\mathrm{red}}
  &:=
  \{(h,h')\in H_0^2:
      \varepsilon(h)=\varepsilon(h')
      \text{ and } \rho_1(\gamma h)\ne\rho_1(\gamma h')\},\\
  \mathcal O_{\mathrm{rec}}
  &:=
  \{(x,y)\in H_s^2:
      \rho_s(x)=\rho_s(y)
      \text{ and } \rho_t(\theta x)\ne\rho_t(\theta y)\},\\
  \mathcal O_{\mathrm{val}}
  &:=
  \{r\in R_s:
      d_t(\bar\theta r)\ne u(d_s r)\}.
\end{align*}
Then
\begin{align*}
  \operatorname{EventRecorded}(\varepsilon,\gamma,\rho_1)
  &\iff \mathcal O_{\mathrm{ev}}=\emptyset,\\
  \operatorname{RecordStable}(\rho_s,\theta,\rho_t)
  &\iff \mathcal O_{\mathrm{rec}}=\emptyset,\\
  \operatorname{FactValueStable}(\bar\theta,d_s,d_t,u)
  &\iff \mathcal O_{\mathrm{val}}=\emptyset,\\
  \operatorname{StableFact}
  &\iff
  \mathcal O_{\mathrm{ev}}=\emptyset
  \wedge
  \mathcal O_{\mathrm{rec}}=\emptyset
  \wedge
  \mathcal O_{\mathrm{val}}=\emptyset,
\end{align*}
together with the record-comparison status partition
\[
  \operatorname{ExactlyOneRecordComparisonStatus}
    (\varepsilon,\gamma,\rho_1),
\]
the retention and carrying conditions for the record-event extension
\(\langle\rho_1\circ\gamma,\varepsilon\rangle:H_0\to R_1\times E\),
\begin{align*}
  &\operatorname{RetainsRecord}
      \bigl(\rho_1\circ\gamma,
        \langle\rho_1\circ\gamma,\varepsilon\rangle\bigr),\\
  &\operatorname{CarriesEvent}
      \bigl(\langle\rho_1\circ\gamma,\varepsilon\rangle,
        \varepsilon\bigr),
\end{align*}
and the characterization of idempotent packaging by its image,
\[
  \operatorname{IdempotentPackaging}(C)
  \iff
  C(H)=\operatorname{Fix}(C),
\]
where \(\operatorname{IdempotentPackaging}(C):\Longleftrightarrow\forall h\in H,\ C(C(h))=C(h)\)
and \(\operatorname{Fix}(C)=\{h:C(h)=h\}\). Here, with \(s=\rho_1\circ\gamma\),
\(\operatorname{EventRecorded}(\varepsilon,\gamma,\rho_1):\Longleftrightarrow\exists\,\bar\varepsilon:R_1\to E,\ \bar\varepsilon\circ s=\varepsilon\);
\(\operatorname{RecordStable}(\rho_s,\theta,\rho_t):\Longleftrightarrow\exists\,\theta^\flat:R_s\to R_t,\ \theta^\flat\circ\rho_s=\rho_t\circ\theta\);
\(\operatorname{FactValueStable}(\bar\theta,d_s,d_t,u):\Longleftrightarrow d_t\circ\bar\theta=u\circ d_s\);
and \(\operatorname{StableFact}\) is the conjunction of these three.
The record-comparison status takes one of four values
(\emph{exact}, \emph{redundant}, \emph{ambiguous}, or
\emph{incomparable}) determined by the joint emptiness pattern of
\(\mathcal O_{\mathrm{ev}}\) and \(\mathcal O_{\mathrm{red}}\): exact when
both are empty, redundant when only \(\mathcal O_{\mathrm{ev}}\) is empty,
ambiguous when only \(\mathcal O_{\mathrm{red}}\) is empty, and incomparable
when neither is.
If moreover \(\bar\theta\) is a descended transport,
\(\bar\theta\circ\rho_s=\rho_t\circ\theta\), then fact-value stability is the
raw fact-value equation:
\[
  \operatorname{FactValueStable}(\bar\theta,d_s,d_t,u)
  \iff
  \forall x\in H_s,\ d_t(\rho_t(\theta x))=u(d_s(\rho_s x)).
\]
If moreover \(H_s=H_1\), \(\rho_s=\rho_1\), \(E_s=E\),
\(d_s\circ\rho_1\circ\gamma=\varepsilon\) and
\(\bar\theta\circ\rho_1=\rho_t\circ\theta\), then
\(\operatorname{FactValueStable}(\bar\theta,d_s,d_t,u)\) implies
\(d_t\circ\rho_t\circ\theta\circ\gamma=u\circ\varepsilon\): a stable fact
transports the event it records.
\end{theorem}
\begin{proof}
Assume the displayed maps, the surjectivity of \(s=\rho_1\circ\gamma\),
and the surjectivity of \(\rho_s\).  We prove the conjuncts in
turn.

\emph{(1) Event recorded iff event-record obstruction empty.}
This is the descent criterion of \Cref{sec:apparatus} for the
surjective record signature \(s\) and the readout \(\varepsilon\): the
descended event readout is \(\bar\varepsilon(s(h)):=\varepsilon(h)\).

\emph{(2) Record stable iff record-stability obstruction empty.}
The record-stability equivalence is the same quotient argument for
\(\rho_s\).  The condition \(\mathcal O_{\mathrm{rec}}=\emptyset\) is
exactly the independence needed to define a map
\[
  \theta^{\flat}(\rho_s x):=\rho_t(\theta x),
  \qquad\text{so that}\qquad
  \theta^{\flat}\circ\rho_s=\rho_t\circ\theta.
\]
The fact-value clause below concerns the given map \(\bar\theta\);
when record stability holds, the natural choice of \(\bar\theta\) is
this descended transport \(\theta^{\flat}\). For the last clause, if
\(\bar\theta\circ\rho_s=\rho_t\circ\theta\), then
\(d_t(\rho_t(\theta x))=d_t(\bar\theta(\rho_s x))\), so the raw equation at
\(x\) is the fact-value equation at \(\rho_s x\); since \(\rho_s\) is
surjective, every \(r\in R_s\) is of this form.

\emph{(3) Fact-value stable iff the fact-value defect is empty.}
The fact-value defect at \(r\in R_s\) is the failure of the
commutative square
\[
  d_t\circ\bar\theta=u\circ d_s.
\]
Equivalently,
\[
  \mathcal O_{\mathrm{val}}=\emptyset
  \quad\Longleftrightarrow\quad
  \forall r\in R_s,\ d_t(\bar\theta r)=u(d_s r),
\]
which is precisely \(\operatorname{FactValueStable}(\bar\theta,d_s,d_t,u)\).

\emph{(4) Stable fact iff the three obstructions are empty.}
By definition, a stable fact requires event recording, record
stability, and fact-value stability.  Combining (1)--(3) gives
\[
  \operatorname{StableFact}
  \iff
  \mathcal O_{\mathrm{ev}}=\emptyset
  \wedge
  \mathcal O_{\mathrm{rec}}=\emptyset
  \wedge
  \mathcal O_{\mathrm{val}}=\emptyset.
\]

\emph{(5) Exactly-one record-comparison status.}
The four statuses are decided by the joint
\((\mathcal O_{\mathrm{ev}},\mathcal O_{\mathrm{red}})\) emptiness
pattern:
\[
\begin{array}{c|c|c}
\text{status} & \mathcal O_{\mathrm{ev}}=\emptyset
  & \mathcal O_{\mathrm{red}}=\emptyset\\
\hline
\text{exact record} & \text{yes} & \text{yes}\\
\text{redundant record} & \text{yes} & \text{no}\\
\text{ambiguous record} & \text{no} & \text{yes}\\
\text{incomparable record} & \text{no} & \text{no}
\end{array}
\]
Every configuration of the two Boolean defect tests falls into one
row, and no configuration falls into two rows.  Existence is by case
enumeration; exclusivity is by definitional disjointness of the four
status predicates.

\emph{(6) Retains record.}
Define the canonical record-event extension
\[
  \langle\rho_1\circ\gamma,\varepsilon\rangle:
  H_0\longrightarrow R_1\times E,
  \qquad
  h\longmapsto(\rho_1(\gamma h),\varepsilon(h)).
\]
Projection to the first coordinate recovers the record signature:
\[
  \operatorname{pr}_{R_1}
  \circ\langle\rho_1\circ\gamma,\varepsilon\rangle
  =
  \rho_1\circ\gamma.
\]

\emph{(7) Carries event.}
Projection to the second coordinate recovers the event readout:
\[
  \operatorname{pr}_{E}
  \circ\langle\rho_1\circ\gamma,\varepsilon\rangle
  =
  \varepsilon.
\]

\emph{(8) Idempotent packaging is image equal to fixed locus.}
Always \(\operatorname{Fix}(C)\subseteq C(H)\), since \(h=C(h)\) for a fixed
point. If \(C(C(h))=C(h)\) for every \(h\), then each \(C(h)\) is fixed, so
\(C(H)\subseteq\operatorname{Fix}(C)\). Conversely, if
\(C(H)\subseteq\operatorname{Fix}(C)\), then \(C(h)\) is fixed for every
\(h\), that is, \(C(C(h))=C(h)\).

\emph{(9) A stable fact transports its event.} Under the additional
identifications, for every \(h\in H_0\),
\(d_t\rho_t\theta\gamma h=d_t\bar\theta\rho_1\gamma h=u\,d_s\rho_1\gamma h
=u\,\varepsilon h\), using the descended transport, fact-value stability and
the decoding identity in turn.
\end{proof}

\paragraph{Nonclaims.} The theorem does not determine which candidates count as evidence and
does not identify the fact with the record.

\subsection*{Layer instantiations}\label{layer-inst:F20}

\begingroup\small
\begin{description}
\item[\emph{{FAIR} scientific-data principles (research
methodology).}] \textit{Analogy.}
The FAIR Guiding Principles --- Findable, Accessible,
Interoperable, Reusable --- are guidelines for the findability and
reuse of scientific data.  Read through \Fref{F20}, data form stable
facts when generation events are recorded as provenance
metadata, records persist under storage / migration transport via
persistent identifiers and access protocols, and derived analyses
preserve fact-value consistency through declared processing
pipelines.  Each obstruction (unrecorded provenance, broken DOI,
undocumented pipeline change) is an \Fref{F20}-style stable-fact failure
\citep{WilkinsonDumontier2016FAIRPrinciples}.

\end{description}
\endgroup


\subsection{Reflexive Nonclosure [\texorpdfstring{\Fref{F21}}{F21}]}\label{sec:status-records-coherence:reflexive-nonclosure}

\paragraph{Reading the statement.} The symbols are those of \Cref{thm:F21} below.
\(\operatorname{CompleteSelfSoundness}_{\pi}(c)\) says that the claim
\(c\) asserts the complete soundness of the policy at its own layer; a
scoped lower audit is an audit of a more limited claim below that
assertion. \(\operatorname{InScope}_{\pi}\) and
\(\operatorname{SameLayerCycle}_{\pi}\) mark the claims the policy
covers and those that depend on themselves at the same layer. Rules 2 and 3 fix the statuses of outside-scope and in-scope cyclic claims, and the rule on \(\mu\)
says that the meta layer does not accept a meta-claim about its own
soundness. Rule 4 is the premise through which same-layer circularity enters;
the blocking of self-certification is derived from rules 2--4, each needed by
the sharpness clause, and the meta-layer clause is the contrapositive of the
rule on \(\mu\). The scoped-lower-audit clause follows from the blocking of self-certification.

\begin{theorem}[Reflexive Nonclosure]\label{thm:F21}
Let \(\mathcal C\) be a space of claims and let \(\pi\) be a same-layer
audit policy on \(\mathcal C\): predicates
\(\operatorname{InScope}_{\pi}\),
\(\operatorname{CompleteSelfSoundness}_{\pi}\),
\(\operatorname{ScopedLowerAudit}_{\pi}\) and
\(\operatorname{SameLayerCycle}_{\pi}\) on \(\mathcal C\), and a status
map \(\operatorname{st}_{\pi}\) from \(\mathcal C\) to a set of statuses
containing three distinct statuses \(\mathsf{accepted}\),
\(\mathsf{outsideScope}\) and \(\mathsf{undefinedCircular}\), subject to rules 2--4 below (rule 1 describes the intended policies and is not assumed):
\begin{enumerate}
\item (intended, not assumed) a scoped lower audit is not a complete self-soundness claim;
\item a claim outside the scope receives \(\mathsf{outsideScope}\);
\item a claim in scope with a same-layer cycle receives
  \(\mathsf{undefinedCircular}\);
\item a complete self-soundness claim in scope has a same-layer cycle.
\end{enumerate}
Let \(\mu\) be a meta-layer audit into a space \(\mathcal M\) of
meta-claims: a relation \(\operatorname{Audits}_{\mu}(m,c)\) and
predicates \(\operatorname{Accepted}_{\mu}\) and
\(\operatorname{MetaSelfSound}_{\mu}\) on \(\mathcal M\), such that \(\operatorname{MetaSelfSound}_{\mu}(m)\Rightarrow\neg\operatorname{Accepted}_{\mu}(m)\) for every \(m\in\mathcal M\): no
meta-claim asserting its own soundness is accepted by the meta layer.
Write
\begin{align*}
  \operatorname{SelfCertificationBlocked}_{\pi}(c)
  &:\iff
  \bigl(\operatorname{CompleteSelfSoundness}_{\pi}(c)
    \Rightarrow \operatorname{st}_{\pi}(c)\ne\mathsf{accepted}\bigr),\\
  \operatorname{AcceptedScopedLowerAudit}_{\pi}(c)
  &:\iff
  \operatorname{ScopedLowerAudit}_{\pi}(c)
    \wedge\operatorname{st}_{\pi}(c)=\mathsf{accepted},\\
  \operatorname{HasMetaLayerCertificate}_{\mu}(c)
  &:\iff
  \exists m\in\mathcal M,\
    \operatorname{Audits}_{\mu}(m,c)\wedge\operatorname{Accepted}_{\mu}(m).
\end{align*}
Then for every \(c\in\mathcal C\),
\begin{align*}
  &\operatorname{SelfCertificationBlocked}_{\pi}(c),\\
  &\operatorname{AcceptedScopedLowerAudit}_{\pi}(c)
    \Longrightarrow
    \neg\operatorname{CompleteSelfSoundness}_{\pi}(c),\\
  &\operatorname{HasMetaLayerCertificate}_{\mu}(c)\\
  &\qquad{}\Longrightarrow
    \exists m\in\mathcal M,\,
      \operatorname{Audits}_{\mu}(m,c)
      \wedge\operatorname{Accepted}_{\mu}(m)
      \wedge
      \neg\operatorname{MetaSelfSound}_{\mu}(m),\\
  &\operatorname{CompleteSelfSoundness}_{\pi}(c)\\
  &\qquad{}\Longrightarrow
    \bigl(\neg\operatorname{InScope}_{\pi}(c)
      \wedge\operatorname{st}_{\pi}(c)=\mathsf{outsideScope}\bigr)\\
  &\hspace{5em}\vee
    \bigl(\operatorname{InScope}_{\pi}(c)
      \wedge\operatorname{SameLayerCycle}_{\pi}(c)
      \wedge\operatorname{st}_{\pi}(c)=\mathsf{undefinedCircular}\bigr).
\end{align*}
Rules~2, 3 and~4 cannot be dropped: there is a policy on two claims (without rule 4) satisfying
rules 1--3 in which \(\operatorname{SelfCertificationBlocked}_{\pi}\) fails,
a policy on one claim (without rule 2) satisfying rules 1, 3 and 4 in which it fails, and a
policy on one claim (without rule 3) satisfying rules 1, 2 and 4 in which it fails. Rule~1 is not assumed: the second line follows from the first, since a claim with status \(\mathsf{accepted}\) is not a complete self-soundness claim.
\end{theorem}
\begin{proof}
Assume the same-layer policy \(\pi\) and the meta-layer audit
\(\mu\), and fix \(c\in\mathcal C\).  Suppose
\(\operatorname{CompleteSelfSoundness}_{\pi}(c)\).  There are two
cases.  If \(c\) is outside the scope of \(\pi\), rule~2 gives
\(\operatorname{st}_{\pi}(c)=\mathsf{outsideScope}\).  If \(c\) is in
scope, rule~4 gives a same-layer cycle, and rule~3 then gives
\(\operatorname{st}_{\pi}(c)=\mathsf{undefinedCircular}\).  This proves
the fourth line, and in both cases
\(\operatorname{st}_{\pi}(c)\ne\mathsf{accepted}\), which is
\(\operatorname{SelfCertificationBlocked}_{\pi}(c)\): a claim certifying
its own complete soundness at the same layer is never accepted.

If \(\operatorname{AcceptedScopedLowerAudit}_{\pi}(c)\) holds, then
\(\operatorname{st}_{\pi}(c)=\mathsf{accepted}\), so
\(\neg\operatorname{CompleteSelfSoundness}_{\pi}(c)\) by the first line: the audit is
accepted for its scoped lower claim, not for the full assertion of its
own soundness.

If \(\operatorname{HasMetaLayerCertificate}_{\mu}(c)\) holds, choose
\(m\) with \(\operatorname{Audits}_{\mu}(m,c)\) and
\(\operatorname{Accepted}_{\mu}(m)\).  If
\(\operatorname{MetaSelfSound}_{\mu}(m)\) held, the rule on \(\mu\)
would give \(\neg\operatorname{Accepted}_{\mu}(m)\); hence
\(\neg\operatorname{MetaSelfSound}_{\mu}(m)\).  The certificate is
supplied by a meta-claim that does not certify its own soundness, so
the audit has moved to a strict meta-layer.

For the sharpness claim, let \(\mathcal C=\{a,b\}\), put both claims in scope,
let \(b\) be the only complete self-soundness claim and \(a\) the only scoped
lower audit, let no claim have a same-layer cycle, and let
\(\operatorname{st}_{\pi}\) assign \(\mathsf{accepted}\) to both. Rule~1 holds
because \(a\) is not a complete self-soundness claim, and rules~2 and~3 hold
vacuously. Rule~4 fails at \(b\), and \(b\) is a complete self-soundness claim
with status \(\mathsf{accepted}\), so
\(\operatorname{SelfCertificationBlocked}_{\pi}(b)\) fails. For rule~2, let
\(\mathcal C=\{c\}\) with \(c\) a complete self-soundness claim outside the
scope of \(\pi\), with no scoped lower audit, no same-layer cycle, and status \(\mathsf{accepted}\).
Rule~1 holds because there is no scoped lower audit, rules~3 and~4 because
\(c\) is not in scope, rule~2 fails at \(c\), and
\(\operatorname{SelfCertificationBlocked}_{\pi}(c)\) fails. For rule~3, let
\(\mathcal C=\{c\}\) with \(c\) a complete self-soundness claim in scope, with a
same-layer cycle, no scoped lower audit, and status \(\mathsf{accepted}\).
Rule~1 holds because there is no scoped lower audit, rule~2 holds vacuously,
rule~4 holds because \(c\) has a same-layer cycle, rule~3 fails at \(c\), and
\(\operatorname{SelfCertificationBlocked}_{\pi}(c)\) fails.
\end{proof}

\paragraph{Nonclaims.} The theorem does not forbid scoped self-audit; it forbids unrestricted
same-layer self-closure.  A proof-theoretic analogue is G\"odel's second
incompleteness theorem \citep{Godel1931Unentscheidbare}, where unprovability
of self-consistency is derived from the arithmetic of the theory; here
non-acceptance of complete self-soundness follows from rules 2--4 and the
distinctness of the three statuses, which
the policy supplies. The modal-logic reformulation in
terms of a provability operator is developed in
\citet{Boolos1993LogicProvability}.

\subsection*{Layer instantiations}\label{layer-inst:F21}

\begingroup\small
\begin{description}
\item[\emph{{ISO}/{IEC} 17025 laboratory accreditation
(industrial metrology).}] \textit{Analogy.}
Under ISO/IEC 17025, a testing or calibration laboratory
cannot self-certify its measurement competence: an
accredited laboratory must be assessed by a recognised
accreditation body (an ILAC member such as A2LA, UKAS, or
DAkkS), which acts as the meta-layer.  Internal
quality-management records are a scoped lower audit,
sufficient
for internal corrective action but not for legal or
commercial recognition of competence; the meta-layer
accreditation certificate is the legitimate reflexive
audit witness \citep{ISOIEC17025_2017}.

\end{description}
\endgroup


\subsection{Interface Mediation [\texorpdfstring{\Fref{F22}}{F22}]}\label{sec:status-records-coherence:interface-mediation}

\paragraph{Reading the statement.} The symbols are those of \Cref{thm:F22} below.
\(\operatorname{InterfaceMediated}(q,\beta,U,q')\) says that the target
class \(q'(U(x,y))\) depends only on the source class \(q(x)\) and the
boundary value \(\beta(y)\). \(\operatorname{BoundaryLeakageAbsent}\)
says that, for each fixed \(x\), two inputs \(y,y'\) with the same
boundary value give the same target class;
\(\operatorname{InternalInsufficiencyAbsent}\) says that, for each
fixed \(y\), two states \(x,x'\) with the same source class give the
same target class. \(\operatorname{ExternallyInert}(\bar U)\) says that
\(\bar U(a,b)\) does not depend on the boundary argument \(b\), and
\(\operatorname{ActiveMediatedInfluence}(\bar U)\) says that it does
for some \(a\).

\begin{theorem}[Interface Mediation]\label{thm:F22}
Let \(q:X\to Q\) be a source quotient, let
\(\beta:Y\to B\) be a boundary readout, let
\(U:X\times Y\to X'\) be an interface update, and let
\(q':X'\to Q'\) be the target quotient.  Assume \(q\) and \(\beta\) are surjective.  Define
\[
  \mathcal O_U
  :=
  \{((x,y),(x',y')) :
    q(x)=q(x'),\ \beta(y)=\beta(y'),\
    q'(U(x,y))\ne q'(U(x',y'))\}.
\]
Then
\begin{align*}
  \operatorname{InterfaceMediated}(q,\beta,U,q')
  &\iff \mathcal O_U=\emptyset\\
  &\iff
  \exists!\,\bar U:Q\times B\to Q'\ \text{with}\
    \bar U(q(x),\beta(y))=q'(U(x,y))\ \forall x,y,\\
  \mathcal O_U=\emptyset
  &\iff
  \operatorname{BoundaryLeakageAbsent}\wedge
  \operatorname{InternalInsufficiencyAbsent},
\end{align*}
and, by the definitions of active mediated influence and external inertness,
for every \(\bar U:Q\times B\to Q'\), in particular the descended \(\bar U\) of the
second line,
\[
  \operatorname{ActiveMediatedInfluence}(\bar U)
  \iff \neg\operatorname{ExternallyInert}(\bar U).
\]
\end{theorem}
\begin{proof}
Assume the source quotient \(q\), boundary readout \(\beta\),
interface update \(U\), and target quotient \(q'\).  If
\(\operatorname{InterfaceMediated}(q,\beta,U,q')\) holds, then
changing representatives inside the same \(q\)-class and
\(\beta\)-class cannot change the target \(q'\)-class; hence
\(\mathcal O_U=\emptyset\).

Conversely, suppose \(\mathcal O_U=\emptyset\).  Define
\[
  \bar U(q(x),\beta(y)):=q'(U(x,y)).
\]
If \(q(x)=q(x')\) and \(\beta(y)=\beta(y')\), the emptiness of
\(\mathcal O_U\) gives
\[
  q'(U(x,y))=q'(U(x',y')),
\]
so \(\bar U\) is independent of representatives.  Thus the
interface update descends through \(Q\times B\).  The absence of
\(\mathcal O_U\) also rules out the two possible mediation failures:
boundary-equivalent inputs cannot leak distinct target classes, and
source-equivalent internal states cannot create an undeclared target
difference.  Conversely, assume both absences, and let \(q(x)=q(x')\) and
\(\beta(y)=\beta(y')\).  Boundary leakage is absent at the fixed state
\(x\), and internal insufficiency is absent at the fixed input \(y'\), so
\[
  q'(U(x,y))=q'(U(x,y'))=q'(U(x',y')),
\]
and \(\mathcal O_U=\emptyset\).  Finally, \(\bar U\) has active mediated influence
exactly when it is not constant in the external argument, which is
the negation of external inertness.
The descended map \(\bar U\) is unique because \(q\times\beta\) is surjective.
\end{proof}

\paragraph{Nonclaims.} The theorem does not determine the interface's internal structure and
does not apply to undeclared interfaces.

\subsection*{Layer instantiations}\label{layer-inst:F22}

\begingroup\small
\begin{description}
\item[\emph{Abstract-data-type encapsulation (software
engineering).}] \textit{Special case.}
In an abstract data type, the carrier \(X\) is the set of
concrete representations satisfying the representation invariant (e.g., a stack realised as a
linked list, a dynamic array, or a deque); the source
quotient \(q\) is the abstraction function, which sends each
representation to the abstract value it represents. Take \(Y\) to be the set of public method calls with concretely encoded arguments, \(\beta\) the map sending a call to its method name and the abstract values of its arguments, \(U(x,y)\) the implementation's post-state and return value, and \(q'\) the abstraction of the post-state paired with the abstract return value. Require argument representations to satisfy their invariants and \(U\) to be total with invariant-satisfying post-states, including declared error outcomes. Take \(Q=q(X)\), \(B=\beta(Y)\) and \(Q'=q'(X')\).  Encapsulation is the
requirement that two representations in the same source
class respond identically (modulo the target quotient)
to the same public-method call; private-state exposure
through the public interface, where two representations of one abstract
value answer the same public call with different abstract results, is the
internal-insufficiency failure mode; boundary leakage would be
two calls at a fixed concrete representation, with the same
boundary value \(\beta(y)=\beta(y')\), that give different
abstract results \citep{LiskovZilles1974ADT}.

\item[\emph{Black-box security reduction
(theoretical cryptography).}] \textit{Analogy.}
A standard security proof shows that any adversary
breaking a scheme can be turned into an algorithm
solving an underlying hard problem; in the black-box
reduction discipline, the reduction accesses the
adversary only through oracle queries.  The carrier is
the adversary's internal state, the source quotient
collapses adversaries with the same query/response
distribution, the boundary is the oracle transcript,
and the update is the adversary's next step.  The
reduction's success thereby depends only on the
boundary I/O profile, not on the adversary's internal
code.  The random-oracle methodology, in which the hash
function is itself modelled as a uniformly random
oracle, is one influential setting within this
discipline \citep{BellareRogaway1993RandomOracles}.

\end{description}
\endgroup


\subsection{Probability as Stable Unresolved Fiber [\texorpdfstring{\Fref{F23}}{F23}]}\label{sec:status-records-coherence:probability}

\paragraph{Reading the statement.} The symbols are those of \Cref{thm:F23} below.
\(\operatorname{Unresolved}(q,r)\) says that some \(q\)-fiber contains
two readout values. \(\operatorname{Stable}\) says that each ledger has its
support contained in its own fiber, at the source and at the target, and that
the two declared maps \(\operatorname{push}\) and \(\operatorname{mix}\),
the transported ledger and the declared mixture of target ledgers, agree on
every source class; \(\operatorname{push}\) reads the source ledger through
\(\Phi\), and \(\operatorname{mix}\) reads the ledger of the continued class
through \(\Psi\).
\(\operatorname{StableProbability}(q,r)\) says that the readout is
unresolved and the declared ledgers satisfy \(\operatorname{Stable}\).
\(\operatorname{Stable}\) constrains \(\lambda\), \(\lambda'\), their supports and
the declared maps \(\Phi\), \(\Psi\), \(\bar\kappa\); it does not involve the
weight assignment \(m\), so the values \(\operatorname{Pr}_{q,r}\) enter only
clauses (ii) and (iv)--(vi). With \(S\) and \(S'\) empty and \(\mathcal T\) a
one-point set, \(\operatorname{Stable}\) holds, and by (iii)
\(\operatorname{StableProbability}(q,r)\) is then equivalent to
\(\operatorname{Unresolved}(q,r)\).  Clause (iii) unfolds the definition
of \(\operatorname{StableProbability}\); no normalization or additivity
of the weights is assumed beyond what clauses (iv)--(vi) state. The relation \(S\) is declared separately from the weights: it constrains \(m(\lambda(a))\) only through the hypothesis of clause (v). Under \(\operatorname{Stable}\), a ledger shared by two classes has empty support, since its support lies in both fibers; clause (vi), which does not use \(S\), allows such sharing.

\begin{theorem}[Probability as Stable Unresolved Fiber]\label{thm:F23}
Let \(q:H\to Q\) be a surjective quotient of histories, let
\(\lambda:Q\to L\) be a ledger map with support relation
\(S\subseteq L\times H\), let
\(m:L\to (H\to\mathrm{Prop})\to W\) be a weight assignment, and let
\(r:H\to Z\) be a readout.  Let a declared continuation carry a target
quotient \(q':H'\to Q'\) with ledger \(\lambda':Q'\to L'\) and support
relation \(S'\subseteq L'\times H'\), and let
\(\bar\kappa:Q\to Q'\) be the class-level continuation, let \(\mathcal T\) be a
set with a declared ledger transport \(\Phi:L\to\mathcal T\) and target
mixture \(\Psi:L'\to\mathcal T\), and put
\(\operatorname{push}:=\Phi\circ\lambda\) and
\(\operatorname{mix}:=\Psi\circ\lambda'\circ\bar\kappa\); only their pointwise
equality in \(\operatorname{Stable}\) is used. The continuation data enter only
through \(\operatorname{Stable}\), which enters clause (iii), the last sentence
of clause (v), and the final equivalence of clause (vi). Apart from these uses
of \(\operatorname{Stable}\), clauses (i), (ii) and (iv)--(vi) depend only on
\(q\), \(\lambda\), \(S\), \(m\), \(r\) and their clause-specific hypotheses.  For \(a\in Q\) and \(v\in Z\), set
\[
  \operatorname{Pr}_{q,r}(a,v)
  :=
  m(\lambda(a))\bigl(\{h\in H:q(h)=a\ \text{and}\ r(h)=v\}\bigr),
\]
and write
\begin{align*}
  \operatorname{Stable}
  :\iff{}&
  \bigl(\forall a\in Q,\forall h\in H,\
    S(\lambda(a),h)\Rightarrow q(h)=a\bigr)\\
  &{}\wedge
  \bigl(\forall a'\in Q',\forall h'\in H',\
    S'(\lambda'(a'),h')\Rightarrow q'(h')=a'\bigr)\\
  &{}\wedge
  \bigl(\forall a\in Q,\
    \operatorname{push}(a)=\operatorname{mix}(a)\bigr),\\
  \operatorname{Unresolved}(q,r)
  :\iff{}&
  \exists h,h'\in H,\ q(h)=q(h')\ \text{and}\ r(h)\ne r(h'),\\
  \operatorname{StableProbability}(q,r)
  :\iff{}&
  \operatorname{Unresolved}(q,r)
  \wedge
  \operatorname{Stable}.
\end{align*}
Then:
\begin{enumerate}[label=(\roman*)]
\item \(r\) descends through \(q\) if and only if
  \(q(h)=q(h')\Rightarrow r(h)=r(h')\) for all \(h,h'\in H\), and exactly
  one of ``\(r\) descends through \(q\)'' and \(\operatorname{Unresolved}(q,r)\)
  holds;
\item if \(r=\bar r\circ q\), then for every \(a\in Q\) and \(v\in Z\) the
  event \(\{h:q(h)=a\ \text{and}\ r(h)=v\}\) is the fiber \(q^{-1}(a)\) when
  \(v=\bar r(a)\) and is empty otherwise, so
  \(\operatorname{Pr}_{q,r}(a,v)=m(\lambda(a))(q^{-1}(a))\) for
  \(v=\bar r(a)\) and \(\operatorname{Pr}_{q,r}(a,v)=m(\lambda(a))(\emptyset)\)
  for \(v\ne\bar r(a)\);
\item \(\operatorname{StableProbability}(q,r)\) holds if and only if
  \(\operatorname{Stable}\) holds and \(r\) does not descend through \(q\);
\item if the weights form a commutative monoid, \(Z\) is finite and each
  \(m(\lambda(a))\) vanishes on the empty event and is additive on disjoint
  events, then
  \(\sum_{v\in Z}\operatorname{Pr}_{q,r}(a,v)=m(\lambda(a))(q^{-1}(a))\) for
  every \(a\in Q\): the readout distributes the weight of the fiber;
\item if the support conjunct of \(\operatorname{Stable}\) holds at \(a\), that
  is, \(S(\lambda(a),h)\Rightarrow q(h)=a\), and there are an operation \(\oplus\) on \(W\) and
  an element \(0\in W\) (no monoid laws are needed) with
  \(m(\lambda(a))(A\cup B)=m(\lambda(a))(A)\oplus m(\lambda(a))(B)\) for disjoint
  events \(A,B\) and \(m(\lambda(a))(E)=0\) for every event \(E\) disjoint from
  \(\{h:S(\lambda(a),h)\}\), then
  \(\operatorname{Pr}_{q,r}(a,v)=m(\lambda(a))(\{h:r(h)=v\})\) for every
  \(v\in Z\): the ledger of a class is concentrated on its fiber; in
  particular, under \(\operatorname{Stable}\) the support hypothesis holds at
  every \(a\in Q\), so this conclusion holds at every \(a\) at which the two
  weight conditions hold;
\item if \(W=\mathbb R\), \(H\) is finite and for each \(a\in Q\) there are weights
  \(w_a:q^{-1}(a)\to(0,1]\) with \(\sum_{h\in q^{-1}(a)}w_a(h)=1\) and
  \(m(\lambda(a))(E)=\sum_{h\in E}w_a(h)\) for every event \(E\subseteq q^{-1}(a)\)
  (no condition is imposed on other events, so distinct classes may share a
  ledger), then \(r\) descends through \(q\) if and only if every
  \(\operatorname{Pr}_{q,r}(a,\cdot)\) is a point mass, that is,
  \(\operatorname{Pr}_{q,r}(a,v)=1\) for some \(v\); hence
  \(\operatorname{Unresolved}(q,r)\) holds exactly when some
  \(\operatorname{Pr}_{q,r}(a,\cdot)\) charges two values; consequently,
  under these weights, \(\operatorname{StableProbability}(q,r)\) holds if and
  only if \(\operatorname{Stable}\) holds and some
  \(\operatorname{Pr}_{q,r}(a,\cdot)\) charges two values.
\end{enumerate}
\end{theorem}
\begin{proof}
Assume the surjective quotient \(q\), the ledgers, supports and
weights, the continuation data, and the readout \(r\).  (i) The first equivalence is the descent criterion of
\Cref{sec:apparatus}, since \(q\) is surjective.
\(\operatorname{Unresolved}(q,r)\) is the negation of the constancy
condition, so exactly one of the two holds.

(ii) Let \(r=\bar r\circ q\) and \(a\in Q\).  Every \(h\) with \(q(h)=a\)
has \(r(h)=\bar r(a)\).  So the event \(\{h:q(h)=a\ \text{and}\ r(h)=v\}\)
is the whole fiber \(q^{-1}(a)\) when \(v=\bar r(a)\) and is empty
otherwise, and the formula for \(\operatorname{Pr}_{q,r}\) follows from its
definition: the readout is deterministic on each fiber.

(iii) \(\operatorname{StableProbability}(q,r)\) is
\(\operatorname{Unresolved}(q,r)\wedge\operatorname{Stable}\), and by (i)
\(\operatorname{Unresolved}(q,r)\) holds exactly when \(r\) does not descend
through \(q\).
Unresolvedness alone is not probability: a stable probability also requires
that each ledger have its support contained in its own fiber, at the source and at the
target, and that transporting the ledger of each source class give the
declared mixture of target ledgers, which is \(\operatorname{Stable}\).

(iv) For fixed \(a\), the events \(\{h:q(h)=a\ \text{and}\ r(h)=v\}\),
\(v\in Z\), are pairwise disjoint and their union is \(q^{-1}(a)\). Finite
additivity of \(m(\lambda(a))\) gives the stated sum.

(v) The event \(\{h:r(h)=v\}\) is the disjoint union of
\(\{h:q(h)=a\ \text{and}\ r(h)=v\}\) and \(\{h:q(h)\ne a\ \text{and}\ r(h)=v\}\).
By the support conjunct, the second is disjoint from
\(\{h:S(\lambda(a),h)\}\), so \(m(\lambda(a))\) vanishes on it; so does the
empty event. Write \(E_{a,v}=\{h:q(h)=a\ \text{and}\ r(h)=v\}\) and \(\oplus\)
for the addition of weights. Additivity on \(E_{a,v}\) and \(\emptyset\) gives
\(\operatorname{Pr}_{q,r}(a,v)=m(\lambda(a))(E_{a,v})\oplus 0\), and additivity
on \(E_{a,v}\) and \(\{h:q(h)\ne a\ \text{and}\ r(h)=v\}\) gives
\(m(\lambda(a))(\{h:r(h)=v\})=m(\lambda(a))(E_{a,v})\oplus 0\), where \(0\) is
the common value on events disjoint from the support. Hence
\(m(\lambda(a))(\{h:r(h)=v\})=\operatorname{Pr}_{q,r}(a,v)\).

(vi) The event \(\{h:q(h)=a,\ r(h)=v\}\) lies in \(q^{-1}(a)\), and each \(h\in
q^{-1}(a)\) has \(w_a(h)>0\), so \(\operatorname{Pr}_{q,r}(a,v)=\sum_{h\in q^{-1}(a),\,r(h)=v}w_a(h)\),
these values sum to \(1\) over \(v\), and \(\operatorname{Pr}_{q,r}(a,v)>0\)
exactly when some \(h\in q^{-1}(a)\) has \(r(h)=v\). Hence
\(\operatorname{Pr}_{q,r}(a,\cdot)\) is a point mass exactly when \(r\) takes
one value on \(q^{-1}(a)\), and by (i) this holds for every \(a\) exactly when
\(r\) descends through \(q\).
\end{proof}

\paragraph{Nonclaims.} The theorem does not determine numerical probability values and does
not apply to resolved fibers.  The measure-theoretic foundation of
probability is \citet{Kolmogorov1933Grundbegriffe}; the branch /
decoherent-histories interpretation that comes closest to a stable
unresolved-fiber reading appears in \citet{Griffiths1984ConsistentHistories}.

\subsection*{Layer instantiations}\label{layer-inst:F23}

\begingroup\small
\begin{description}
\item[\emph{Design-based survey sampling (social
statistics).}] \textit{Analogy.}
For a fixed finite population, the carrier is the set of
possible sample-draw realisations under the chosen
randomisation scheme; the source quotient collapses
draws sharing the same population-level summary; the
readout is the realised sample statistic.  The same
population-level summary supports many distinct sample
realisations.  The sampling design's inclusion-probability
measure is the ledger; the design-induced probability
law on draws is the weight assignment.  The
Horvitz-Thompson estimator and the design-based variance
estimator are the stable probability records derived
from this design measure
\citep{Cochran1977SamplingTechniques}.

\item[\emph{Chain-ladder loss-reserving in non-life
insurance (actuarial science).}] \textit{Analogy.}
The carrier is the space of individual
claim-development histories within an accident year;
the source quotient collapses histories sharing the
same observed cumulative-payment pattern through the
current development cell; the readout is the ultimate
loss, which is unresolved over the observed
development-pattern fibre.  Mack's distribution-free
assumptions on the cumulative payments constitute the
ledger and supply the first-and-second-moment structure
of incremental development factors as the weight
assignment.  The chain-ladder reserve estimator and
Mack's prediction-error formula are the stable
probability records used for reserve setting and
solvency capital \citep{Mack1993ChainLadder}.

\item[\emph{Wright-Fisher stochastic population genetics
(theoretical biology).}] \textit{Analogy.}
The carrier is the set of gamete-sampling realisations
within a generation; the source quotient collapses
realisations to the current-generation allele frequency;
the readout is the realised next-generation allele
count.  Two distinct gamete-sampling outcomes sharing
the same current frequency can produce different
next-generation counts, so the readout is unresolved
over the current-frequency fibre.  The ledger records
the effective population size and any
selection/mutation parameters; the weight assignment is
the Wright-Fisher binomial (or diffusion) transition
kernel they induce.  The resulting transition law is
the stable probability record that descends to the
population-level frequency dynamics
\citep{Wright1931Evolution}.

\item[\emph{Shannon noisy-channel coding (information
theory).}] \textit{Analogy.}
The carrier is the set of (input, output) trajectories
of a noisy memoryless channel; the source quotient
collapses trajectories sharing the same observed output
sequence; the readout is the transmitted input.  The
receiver cannot uniquely recover the input from the
output (unresolved fibre).  The channel input
distribution supplies the ledger; the channel law
supplies the weight assignment; the channel capacity and the noisy-channel coding
theorem are results about this channel that \Fref{F23}
neither states nor implies; \Fref{F23} classifies the input readout as unresolved exactly
when at least one output fibre contains trajectories with
different transmitted inputs
\citep{Shannon1948Mathematical}.

\item[\emph{Weibull life-distribution model in
reliability engineering (engineering reliability).}] \textit{Analogy.}
The carrier is the set of individual component
life trajectories under specified operating conditions;
the source quotient collapses trajectories sharing the
same operating-condition class; the readout is the
realised time-to-failure of a particular component.
Individual times-to-failure are unresolved over the
operating-condition fibre.  The Weibull distribution
function \(F(t)=1-\exp(-(t/\eta)^{\beta})\) is the
standard ledger and supplies a stable two-parameter
hazard model as the weight assignment; fitted shape
\(\beta\) and scale \(\eta\) give the stable
probability record used for population-level
reliability prediction, warranty pricing, and burn-in
design \citep{Weibull1951Statistical}.
\end{description}
\endgroup

\subsection{Causality as Stable Intervention Transport [\texorpdfstring{\Fref{F25}}{F25}]}\label{sec:status-records-coherence:causality}

\paragraph{Reading the statement.} The symbols are those of \Cref{thm:F25} below.
\(\operatorname{ITD}\) and \(\operatorname{PITD}\) say that the target
readout, respectively the kernel outcome, depends only on the context
class and the intervention class. \(\operatorname{Caus}(c,u,v)\) says
that interventions of classes \(u\) and \(v\) at context class \(c\)
have a stable causal effect: both are admissible, the quotient-level
transport gives different results for them, and the residue is stable
in the declared sense; \(\operatorname{ResStable}\) is declared data.
\(\operatorname{ProbCaus}\) is the same for kernel outcomes.

\begin{theorem}[Causality as Stable Intervention Transport]\label{thm:F25}
Let \(q:H\to Q\) be a surjective context quotient, let
\(\alpha:A\to\bar A\) be a surjective intervention abstraction, let
\(\Gamma:H\times A\to H'\) be an intervention transport, let
\(r:H'\to R\) be a target readout, and let
\(K:H\times A\to K_0\) be a kernel outcome (in the intended reading \(K_0\) is a set of probability laws on outcomes and \(K(h,a)\) is the outcome law of intervention \(a\) in context \(h\); the theorem uses \(K\) only as a map).  Let
\(\operatorname{Adm}:Q\times\bar A\to\mathrm{Prop}\) be admissibility,
let
\(\operatorname{ResStable},\operatorname{ProbResStable}:
Q\times\bar A\times\bar A\to\mathrm{Prop}\) be residue-stability
predicates, and let \(\bar\Gamma:Q\times\bar A\to R\) and
\(\bar K:Q\times\bar A\to K_0\) be quotient-level transports.  Define
\[
  \mathcal O_{\Gamma}
  :=
  \{((h,a),(h',a')) :
    q(h)=q(h'),\ \alpha(a)=\alpha(a'),\
    r(\Gamma(h,a))\ne r(\Gamma(h',a'))\},
\]
and \(\mathcal O_K\) in the same way with \(K\) in place of
\(r\circ\Gamma\).  Define intervention-transport descent
\(\operatorname{ITD}\), probabilistic intervention-transport descent
\(\operatorname{PITD}\), and, for \(c\in Q\) and \(u,v\in\bar A\), a stable
causal effect \(\operatorname{Caus}(c,u,v)\) and a stable probabilistic
causal effect \(\operatorname{ProbCaus}(c,u,v)\):
\begin{align*}
  \operatorname{ITD}(q,\alpha,\Gamma,r)
  &:\Longleftrightarrow
  \exists\,g:Q\times\bar A\to R,\
    g(q(h),\alpha(a))=r(\Gamma(h,a))\ \forall h,a,\\
  \operatorname{PITD}(q,\alpha,K)
  &:\Longleftrightarrow
  \exists\,k:Q\times\bar A\to K_0,\
    k(q(h),\alpha(a))=K(h,a)\ \forall h,a,\\
  \operatorname{Caus}(c,u,v)
  &:\Longleftrightarrow
  \operatorname{Adm}(c,u)\wedge\operatorname{Adm}(c,v)
  \wedge\bar\Gamma(c,u)\ne\bar\Gamma(c,v)
  \wedge\operatorname{ResStable}(c,u,v),\\
  \operatorname{ProbCaus}(c,u,v)
  &:\Longleftrightarrow
  \operatorname{Adm}(c,u)\wedge\operatorname{Adm}(c,v)
  \wedge\bar K(c,u)\ne\bar K(c,v)
  \wedge\operatorname{ProbResStable}(c,u,v).
\end{align*}
Then
\begin{align*}
  \operatorname{ITD}(q,\alpha,\Gamma,r)
  &\iff \mathcal O_{\Gamma}=\emptyset
  \iff
  \exists!\,g:Q\times\bar A\to R,\
    g(q(h),\alpha(a))=r(\Gamma(h,a))\ \forall h,a,\\
  \operatorname{PITD}(q,\alpha,K)
  &\iff \mathcal O_{K}=\emptyset
  \iff
  \exists!\,k:Q\times\bar A\to K_0,\
    k(q(h),\alpha(a))=K(h,a)\ \forall h,a.
\end{align*}
Stable causal effects have raw witnesses, each under its own descent
equation. If \(\bar\Gamma\) is a descended transport,
\(\bar\Gamma(q(h),\alpha(a))=r(\Gamma(h,a))\) for all \(h,a\), then
\(\operatorname{Caus}(c,u,v)\) implies that there are \(h\in H\) and
\(a_1,a_2\in A\) with \(q(h)=c\), \(\alpha(a_1)=u\), \(\alpha(a_2)=v\) and
\(r(\Gamma(h,a_1))\ne r(\Gamma(h,a_2))\). If \(\bar K\) is a descended
transport, \(\bar K(q(h),\alpha(a))=K(h,a)\) for all \(h,a\), then
\(\operatorname{ProbCaus}(c,u,v)\) implies the same with \(K\) in place of
\(r\circ\Gamma\).
\end{theorem}
\begin{proof}
Assume the surjective quotients \(q\) and \(\alpha\), the transport
\(\Gamma\), the target readout \(r\), and the kernel outcome \(K\).
The first line is the descent criterion of \Cref{sec:apparatus} for
the surjective source \(q\times\alpha:H\times A\to Q\times\bar A\) and the
readout \((h,a)\mapsto r(\Gamma(h,a))\), with
\(g(q(h),\alpha(a)):=r(\Gamma(h,a))\); the probabilistic statement is
identical with \(K\) in place of \(r\circ\Gamma\).

A causal effect at context class \(c\) compares two intervention
classes \(u\) and \(v\).  It requires both to be admissible at \(c\),
the quotient-level transport to give different results on them (a
nonzero residue), and the residue to be stable; this is the definition of
\(\operatorname{Caus}\), and \(\operatorname{ProbCaus}\) is the same with
\(\bar K\) and \(\operatorname{ProbResStable}\).  A zero residue,
\(\bar\Gamma(c,u)=\bar\Gamma(c,v)\), is a finding of no effect, not a
failure of transport: a transport can descend with every residue zero.

For the last clause, choose by surjectivity \(h\) with \(q(h)=c\) and
\(a_1,a_2\) with \(\alpha(a_1)=u\), \(\alpha(a_2)=v\). Then
\(r(\Gamma(h,a_1))=\bar\Gamma(c,u)\ne\bar\Gamma(c,v)=r(\Gamma(h,a_2))\), using
only the descent equation of \(\bar\Gamma\); the same computation applies to
\(K\) and \(\bar K\), using only the descent equation of \(\bar K\).
Uniqueness of \(g\) and \(k\) holds because \(q\times\alpha:H\times A\to Q\times\bar A\) is surjective.
\end{proof}

\paragraph{Nonclaims.} The theorem does not identify causes with raw interventions and does
not apply to undeclared alternatives.  The intervention-counterfactual
hierarchy that this layer-agnostic formulation echoes is developed
in detail in \citet{Pearl2009Causality}.

\subsection*{Layer instantiations}\label{layer-inst:F25}

\begingroup\small
\begin{description}
\item[\emph{Randomised controlled trial / Fisher
randomisation (experimental design).}] \textit{Analogy.}
The carrier is the population of experimental units
(patients, field plots); the context quotient collapses
units to their pre-randomisation baseline stratum; the
intervention space is the set of treatment arms; the
abstraction collapses operational delivery details to
the trial-arm class; the transport is the trial's
delivery-and-follow-up process; the target readout is
the primary outcome.  Individual outcomes do not
deterministically descend through (baseline stratum,
arm); the \Fref{F25} reading is its probabilistic form:
randomised assignment supplies admissibility, and the
stable residue is the arm-specific outcome distribution
or the within-stratum average treatment effect, secured
by pre-registration, blinding, and intention-to-treat
analysis \citep{Fisher1935DesignExperiments}.

\end{description}
\endgroup


\subsection{Moduli / Contingency [\texorpdfstring{\Fref{F26}}{F26}]}\label{sec:status-records-coherence:moduli}

\paragraph{Reading the statement.} The symbols are those of \Cref{thm:F26} below.
Lawful closures are grouped into moduli by the equivalence \(\sim\), and
\(p\) sends a lawful closure to its modulus. A readout can fail to be a
function of the modulus (it is presentation-dependent), or it can be a
function of the modulus and then either take one value on every lawful
closure (necessary) or take different values on two inequivalent lawful
closures (contingent). The definitions come first; the theorem proves that
these alternatives are exhaustive and exclusive, and that descent through
\(p\) is the same as independence of presentation. The predicate
\(\operatorname{NecessaryBackground}\) of \Cref{thm:F50} is necessity in the
present sense for a background already defined on moduli.

\begin{theorem}[Moduli / Contingency]\label{thm:F26}
Let \(\mathcal C\) be a closure space with lawfulness predicate
\(\operatorname{Lawful}:\mathcal C\to\mathrm{Prop}\), equivalence relation
\(\sim\) on lawful closures (an equivalence relation, by the kernel condition), moduli quotient \(p:\mathcal C\to M\) with
\(p(c)=p(d)\iff c\sim d\) for lawful \(c,d\), and readout
\(r:\mathcal C\to V\). Assume that every \(m\in M\) is \(p(c)\) for some
lawful \(c\). Define, with \(c,d\) ranging over lawful closures,
\begin{align*}
  \operatorname{ClosureObstructed}
  &:\Longleftrightarrow
  \text{no closure is lawful},\\
  \operatorname{NontrivialModuli}
  &:\Longleftrightarrow
  \exists c,d:\ \neg(c\sim d),\\
  \operatorname{ModuliReadout}(r)
  &:\Longleftrightarrow
  \forall c,d:\ c\sim d\Rightarrow r(c)=r(d),\\
  \operatorname{PresentationDependent}(r)
  &:\Longleftrightarrow
  \exists c,d:\ c\sim d\ \wedge\ r(c)\ne r(d),\\
  \operatorname{NecessaryReadout}(r)
  &:\Longleftrightarrow
  \forall c,d:\ r(c)=r(d),\\
  \operatorname{ContingentReadout}(r)
  &:\Longleftrightarrow
  \exists c,d:\ \neg(c\sim d)\ \wedge\ r(c)\ne r(d).
\end{align*}
Then:
\begin{enumerate}[label=(\roman*)]
\item \(\operatorname{ModuliReadout}(r)\) holds if and only if there is
  \(\bar r:M\to V\) with \(\bar r(p(c))=r(c)\) for every lawful \(c\), and
  \(\operatorname{PresentationDependent}(r)\) holds if and only if
  \(\operatorname{ModuliReadout}(r)\) fails;
\item if \(\operatorname{ModuliReadout}(r)\) holds, then exactly one of
  \(\operatorname{NecessaryReadout}(r)\) and
  \(\operatorname{ContingentReadout}(r)\) holds;
\item if \(\operatorname{NontrivialModuli}\) fails, every moduli readout is
  necessary; if \(\operatorname{ClosureObstructed}\) holds, every readout is
  necessary and none is contingent or presentation-dependent, since these
  three predicates quantify over lawful closures only; this part does not use
  the standing assumption. Under the standing assumption,
  \(\operatorname{ClosureObstructed}\) moreover forces \(M=\emptyset\), hence
  \(\mathcal C=\emptyset\), because \(p\) is defined on all of \(\mathcal C\);
\item for a map \(T:\mathcal C\to\mathcal C'\) into a second closure space
  with lawfulness \(\operatorname{Lawful}'\) and equivalence \(\sim'\), taking
  lawful closures to lawful closures, \(T\) respects moduli
  (\(c\sim d\Rightarrow T(c)\sim' T(d)\) for lawful \(c,d\)) if and only if no
  pair of lawful closures \(c\sim d\) has \(\neg(T(c)\sim' T(d))\).
\end{enumerate}
\end{theorem}
\begin{proof}
(i) If \(\bar r\) exists and \(c\sim d\), then \(p(c)=p(d)\), so
\(r(c)=\bar r(p(c))=\bar r(p(d))=r(d)\). Conversely, let \(r\) be a moduli
readout. For \(m\in M\) choose a lawful \(c_m\) with \(p(c_m)=m\) and set
\(\bar r(m):=r(c_m)\). For lawful \(c\), \(p(c_{p(c)})=p(c)\) gives
\(c_{p(c)}\sim c\), so \(\bar r(p(c))=r(c_{p(c)})=r(c)\). The second
statement is the negation of the definition of a moduli readout.

(ii) Let \(r\) be a moduli readout. If \(r\) is necessary, no two lawful
closures have different readouts, so \(r\) is not contingent. If \(r\) is not
necessary, there are lawful \(c,d\) with \(r(c)\ne r(d)\); since \(r\) is a
moduli readout, \(c\not\sim d\), so \(r\) is contingent.

(iii) If \(\operatorname{NontrivialModuli}\) fails, all lawful closures are
equivalent, so a moduli readout takes one value on them. If no closure is
lawful, the universal conditions hold vacuously and the existential ones fail.

(iv) The obstruction pairs are exactly the failures of the implication that
defines respecting moduli.
\end{proof}

\paragraph{Nonclaims.} The theorem does not enumerate a carrier's moduli and does not
collapse stable parameters with contingent choices.

\subsection*{Layer instantiations}\label{layer-inst:F26}

\begingroup\small
\begin{description}
\item[\emph{Moduli space of elliptic curves (algebraic
geometry).}] \textit{Special case.}
The carrier is the set of Weierstrass-form presentations
\(y^2=x^3+ax+b\) over an algebraically closed base field of characteristic other
than 2 and 3; the lawfulness
predicate restricts to non-singular cubics
\((4a^3+27b^2\ne 0)\); the equivalence is isomorphism of
elliptic curves; the moduli quotient sends a
non-singular presentation to its isomorphism class and every singular
presentation to the class of \(y^2=x^3-x\) (\Fref{F26}'s kernel condition is imposed only on lawful
closures). Take \(M\) to be the isomorphism classes of
nonsingular presentations, and extend \(j\) to singular
presentations by \(j(a,b)=1728\). This total readout is a moduli readout: distinct Weierstrass
presentations of isomorphic elliptic curves share the
same j-invariant, and the descended map
\(\bar\jmath:M\to\mathbb A^1\) is a bijection onto the
coarse moduli space.  It is contingent, since non-isomorphic
curves have different j-invariants.  The individual coefficients
\((a,b)\) are presentation-dependent: \((a,b)\) and
\((\lambda^4a,\lambda^6b)\) present isomorphic curves.  The genus, equal
to one on every lawful closure, is a necessary readout
\citep{Mumford1965GIT}.

\item[\emph{Crystallographic space groups and Bravais
lattices (solid-state physics).}] \textit{Special case.}
The carrier is the set of three-dimensional crystalline
arrangements; the lawfulness
predicate requires invariance under a rank-3 lattice of
translations; the equivalence is conjugacy of space groups under
orientation-preserving affine maps (change of origin and
right-handed basis); the moduli quotient sends a crystal to its
space-group label among the 230 space-group types.  Abstract
isomorphism, equivalently conjugacy under all affine maps, merges
the 11 enantiomorphic pairs and gives 219 classes.  The space-group label is a
moduli readout, and a contingent one, since lawful crystals with different
space groups exist (rock salt, \(Fm\bar3m\); \(\alpha\)-quartz,
\(P3_121\)).  The lattice metric (unit-cell lengths
\(a,b,c\), angles \(\alpha,\beta,\gamma\),
atomic-displacement parameters) is presentation-dependent for this
equivalence: lawful crystals with the same space group can have different
lattice metrics \citep{HahnITC2016}.

\item[\emph{Standard Model free parameters as
renormalisation moduli (particle physics).}] \textit{Analogy.}
The carrier is the space of Standard-Model Lagrangians
with the fixed gauge group
SU(3)\(\times\)SU(2)\(\times\)U(1), three-generation
fermion content, and Higgs doublet built into the
lawfulness predicate (renormalisability, gauge
invariance, anomaly cancellation, observed particle
content); the equivalence identifies Lagrangians giving
identical low-energy phenomenology under the
renormalisation group; the moduli quotient sends a
Lagrangian to its renormalised-parameter point.  The 19
parameters of the Standard Model with massless neutrinos (three
gauge couplings, nine charged-fermion masses, four CKM-matrix
parameters, the Higgs mass and vacuum expectation value, and the QCD
vacuum angle), extended by the neutrino-sector parameters, are
contingent moduli readouts; in dimensionless form the fermion
masses are replaced by Yukawa couplings and the Higgs mass by the
quartic self-coupling; predicted phenomenology descends
through the renormalised-parameter point
\citep{Workman2022Review}.

\end{description}
\endgroup


\subsection{Conservation as Orbit Descent [\texorpdfstring{\Fref{F27}}{F27}]}\label{sec:status-records-coherence:conservation}

\paragraph{Reading the statement.}
A declared route or symmetry structure on \(X\) is a set \(\Gamma\) of partial
steps \(\gamma:\operatorname{dom}\gamma\to X\) with
\(\operatorname{dom}\gamma\subseteq X\); a family of maps \(X\to X\) is the case
\(\operatorname{dom}\gamma=X\). A declared relation \(\to\) on \(X\) is the case
\(\Gamma=\{\gamma_{x,x'}:x\to x'\}\) with \(\operatorname{dom}\gamma_{x,x'}=\{x\}\) and
\(\gamma_{x,x'}(x)=x'\); its orbits are the classes of the equivalence generated
by \(\to\), and conservation says \(C(x)=C(x')\) whenever \(x\to x'\). The reaction,
Reidemeister and reduction instances below are declared in this form. The \emph{orbits} are the classes of the
equivalence relation generated by \(x\sim\gamma(x)\) for \(\gamma\in\Gamma\)
and \(x\in\operatorname{dom}\gamma\), and \(q_{\mathrm{orb}}\) is the quotient
map onto them. A quantity \(C:X\to V\) is \emph{conserved} when
\(C(\gamma(x))=C(x)\) for every \(\gamma\in\Gamma\) and
\(x\in\operatorname{dom}\gamma\).

\begin{theorem}[Conservation as Orbit Descent]\label{thm:F27}
Let \(\Gamma\) be a set of partial maps \(\gamma:\operatorname{dom}\gamma\to X\)
with \(\operatorname{dom}\gamma\subseteq X\), let \(q_{\mathrm{orb}}\) be the
quotient map of the equivalence relation generated by \(x\sim\gamma(x)\)
(\(\gamma\in\Gamma\), \(x\in\operatorname{dom}\gamma\)), and let \(C:X\to V\).
Then \(C\) is conserved, \(C(\gamma(x))=C(x)\) for all \(\gamma\in\Gamma\) and
\(x\in\operatorname{dom}\gamma\), if and only if it descends through the orbit
quotient, \(C=\bar C\circ q_{\mathrm{orb}}\) for some \(\bar C\); that is,
\(C(x)=C(x')\) whenever \(x\) and \(x'\) lie in the same orbit.
\end{theorem}
\begin{proof}
If \(C\) descends through the orbit quotient, then, since \(x\) and
\(\gamma(x)\) lie in the same orbit for \(x\in\operatorname{dom}\gamma\), \(C(\gamma(x))=C(x)\), so \(C\) is
conserved.  Conversely, let \(C\) be conserved. The relation
\(C(x)=C(x')\) is an equivalence relation containing every generating pair
\((x,\gamma(x))\), so it contains the generated equivalence: \(C\) is constant
on orbits.  Define the quotient readout by \(\bar C([x])=C(x)\).  Orbit constancy
makes the definition independent of representative, and
\(C=\bar C\circ q_{\mathrm{orb}}\).  Thus conservation is orbit
descent.
\end{proof}

For the conserved quantity to be lawful at the current layer in the sense of \Cref{thm:F18}, \(C\) must also factor through the access quotient \(\accessQ\) of that layer; descent through \(q_{\mathrm{orb}}\) alone does not give this.

\paragraph{Nonclaims.} The theorem does not enumerate conserved quantities and does not
apply to undeclared routes or symmetries.  Noether's theorem
\citep{Noether1918Invariante} produces conserved quantities from continuous
variational symmetries; the theorem above reads each such quantity as a
function on the orbits of the flow.

\subsection*{Layer instantiations}\label{layer-inst:F27}

\begingroup\small
\begin{description}
\item[\emph{Lavoisier conservation of mass in chemical
reactions (chemistry).}] \textit{Special case.}
The carrier is the set of finite multisets of molecules, each molecule a finite multiset of atoms of declared elements \(e\) with masses \(m_e\); the declared route structure
is the set of elementary chemical reactions that
preserve atomic identity; the orbit quotient refines the map sending a
configuration to its multiset of atoms.  Total mass
\(C(x)=\sum_e n_e(x)\,m_e\) is constant across orbits
of chemical rearrangement and so descends through the
orbit quotient.  Lavoisier's stoichiometric balancing
of chemical equations is the founding expression of
this orbit-descent reading \citep{Lavoisier1789Traite}.

\item[\emph{Knot invariants under Reidemeister moves
(combinatorial knot theory).}] \textit{Special case.}
The carrier is the set of planar knot diagrams; the
declared route structure is the Reidemeister moves (R1
twist, R2 poke, R3 slide) together with planar isotopy;
the orbit quotient sends a diagram to its
ambient-isotopy class.  Polynomial invariants such as
the Jones polynomial or the Conway-normalized Alexander polynomial are unchanged under
every Reidemeister move and so descend through the
orbit quotient to functions on knots.  Reidemeister's
theorem certifies that ambient isotopy of knots is
exactly generated by the three planar moves
\citep{Reidemeister1927Knotentheorie}.

\item[\emph{Subject reduction in operational semantics
(programming-language theory).}] \textit{Special case.}
The carrier is the set of typable program expressions
under a fixed typing context \(\Gamma\) and a type
system in which each typable expression has a unique type (as in the
Church-style simply typed \(\lambda\)-calculus); the declared route structure is the small-step
reduction relation \(\to\) with its reachability
preorder.  Subject reduction (type preservation) is the
conservation statement: if
\(\Gamma\vdash e:\tau\) and \(e\to e'\) then
\(\Gamma\vdash e':\tau\).  On the typable carrier every reduction step therefore links
two expressions of the same type, so the type readout is constant on
each class of the equivalence generated by \(\to\) and descends
through this orbit quotient; no subject expansion is needed, because
both ends of each step are typable.  Together with progress, it yields
the Wright-Felleisen formulation of type soundness
\citep{WrightFelleisen1994SyntacticTypeSoundness}.

\item[\emph{Meselson-Stahl semi-conservative {DNA}
replication (molecular biology).}] \textit{Analogy.}
The carrier is the population of DNA duplexes in a
dividing cell; the declared route structure is one
round of replication, in which each parental duplex
yields two daughter duplexes that each contain one
parental strand and one newly synthesised strand; the
orbit quotient sends a daughter duplex to its parental
template-strand identity.  The conserved quantity is
the template-strand sequence content carried by each
daughter molecule, which is constant across the
replication orbit (modulo replication-error
corrections).  The Meselson-Stahl density-gradient
experiment is the canonical demonstration,
distinguishing semi-conservative replication from
conservative and dispersive alternatives
\citep{MeselsonStahl1958Replication}.

\item[\emph{Martingale property under conditional
expectation (probability theory).}] \textit{Special case.}
For a stochastic process \((M_t)_{t\ge 0}\) with \(\mathbb E|M_t|<\infty\) for every \(t\), adapted to a filtration
\((\mathcal F_t)\) and a fixed time \(s\), the carrier is the time set
\([s,\infty)\); the declared route structure is time advance, whose
orbit quotient collapses \([s,\infty)\) to a single class; the
quantity is \(C_s(t)=\mathbb E[M_t\mid\mathcal F_s]\), valued in
integrable random variables modulo almost-sure equality.  The
martingale property is this conservation statement for every \(s\):
\(\mathbb E[M_t\mid\mathcal F_s]=M_s\) almost surely for all
\(t\ge s\), so \(C_s\) is constant along
the orbit, equal to its value \(M_s\) at \(t=s\).  The
\(\mathcal F_s\)-measurability of \(\mathbb E[M_t\mid\mathcal F_s]\)
holds for every integrable process and is not the content.  Optional
stopping, Doob's maximal inequalities, and martingale convergence
are built on this conservation property
\citep{Doob1953Stochastic}.
\end{description}
\endgroup

\subsection{Composability [\texorpdfstring{\Fref{F28}}{F28}]}\label{sec:status-records-coherence:composability}

\paragraph{Reading the statement.} The symbols are those of \Cref{thm:F28} below.
\(\operatorname{COD}\), \(\operatorname{CRD}\) and
\(\operatorname{AD}\) say that the composite outcome, the cross residual
and admissibility depend only on the parts' classes and interface
values \(\kappa(a,b)\). \(\operatorname{Composable}\) is exact
composability, the conjunction of the three. \(\operatorname{Assoc}\)
and \(\operatorname{Id}\) are the coherence equations named in the
display.

\begin{theorem}[Composability]\label{thm:F28}
Let \(X_A,X_B\) be component spaces with quotients
\(q_A:X_A\to Q_A\), \(q_B:X_B\to Q_B\), interface readouts
\(\beta_A:X_A\to I_A\), \(\beta_B:X_B\to I_B\), composition
\(U:X_A\times X_B\to X_{AB}\), composite quotient
\(q_{AB}:X_{AB}\to Q_{AB}\), admissibility predicate
\(\operatorname{Adm}:X_A\times X_B\to\mathrm{Prop}\), and residual
\(\delta:X_A\times X_B\to D\).  Write
\[
  \kappa(a,b):=\bigl((q_A(a),\beta_A(a)),(q_B(b),\beta_B(b))\bigr)
\]
for the part/interface quotient and assume \(\kappa\) is surjective.
For a map \(f\) on \(X_A\times X_B\), let
\[
  \mathcal O(f)
  :=
  \{((a,b),(a',b')) :
     \kappa(a,b)=\kappa(a',b'),\ f(a,b)\ne f(a',b')\},
\]
and set \(\mathcal O_U:=\mathcal O(q_{AB}\circ U)\),
\(\mathcal O_{\delta}:=\mathcal O(\delta)\) and
\(\mathcal O_{\operatorname{Adm}}:=\mathcal O(\operatorname{Adm})\).
Write \(\operatorname{COD}\), \(\operatorname{CRD}\) and
\(\operatorname{AD}\) for descent through \(\kappa\) of the composite
outcome, the cross residual and admissibility, and
\(\operatorname{Composable}\) for exact composability.  Then
\begin{align*}
  \operatorname{COD}
  &\iff
  \mathcal O_U=\emptyset
  \iff
    \exists!\,\bar U:(Q_A\times I_A)\times(Q_B\times I_B)\to Q_{AB}\\
  &\hspace{6em}\text{with}\ \bar U(\kappa(a,b))=q_{AB}(U(a,b)),\\
  \operatorname{CRD}
  &\iff
  \mathcal O_{\delta}=\emptyset,
  \qquad
  \operatorname{AD}
  \iff
  \mathcal O_{\operatorname{Adm}}=\emptyset,\\
  \operatorname{Composable}
  &\iff
  \mathcal O_{\operatorname{Adm}}=\emptyset
  \wedge\mathcal O_U=\emptyset
  \wedge\mathcal O_{\delta}=\emptyset.
\end{align*}
The maps through which the cross residual and admissibility descend are likewise unique.
Associativity and identity are separate coherence conditions: for the
two bracketings \(\ell,\varrho:Y\to Y'\) of a declared quotient-level triple
composite, define \(\operatorname{Assoc}:\Longleftrightarrow\ell=\varrho\), and
for a declared composite \(c:Z\times E\to Z\) and a declared \(e\in E\) define
\(\operatorname{Id}:\Longleftrightarrow\forall z\in Z,\ c(z,e)=z\); then
\(\operatorname{Assoc}\iff\{y:\ell(y)\ne\varrho(y)\}=\emptyset\) and
\(\operatorname{Id}\iff\{z\in Z:c(z,e)\ne z\}=\emptyset\). If moreover \(\ell\circ\kappa_3=q'\circ L\) and \(\varrho\circ\kappa_3=q'\circ P\) for raw bracketings \(L,P:X_3\to X'\), a surjection \(\kappa_3:X_3\to Y\) and a map \(q':X'\to Y'\), then \(\operatorname{Assoc}\iff q'\circ L=q'\circ P\); in particular raw associativity \(L=P\) gives \(\operatorname{Assoc}\). Likewise, if \(c(\kappa_Z x,e)=\kappa_Z(\tilde c(x))\) for a surjection \(\kappa_Z\) and a map \(\tilde c\), then \(\operatorname{Id}\iff\kappa_Z\circ\tilde c=\kappa_Z\).
\end{theorem}
\begin{proof}
Assume the component quotients \(q_A,q_B\), interface readouts
\(\beta_A,\beta_B\), composition \(U\), composite quotient
\(q_{AB}\), admissibility predicate \(\operatorname{Adm}\), residual
\(\delta\), and surjectivity of \(\kappa\).  The composite outcome
descends exactly when it is independent of representatives of the
\(\kappa\)-classes, that is, of the parts' classes and interface values.
Thus, if \(\mathcal O_U=\emptyset\), the definition
\[
  \bar U(\kappa(a,b)):=q_{AB}(U(a,b))
\]
is well-defined, and surjectivity of \(\kappa\) defines it everywhere; two such maps agree on the image of the surjection \(\kappa\), so they are equal;
if \(\mathcal O_U\ne\emptyset\), a pair in \(\mathcal O_U\) is precisely
a split-pair obstruction.  The same argument applied to \(\delta\) and
to \(\operatorname{Adm}\) gives the second line, and exact
composability is the conjunction of the three descent conditions.
Associativity and identity are pointwise equations, so each holds
exactly when its set of failures is empty. For the raw clauses, precompose with the surjections: \(\ell=\varrho\) iff \(\ell\circ\kappa_3=\varrho\circ\kappa_3\) iff \(q'\circ L=q'\circ P\), and \(c(\cdot,e)=\mathrm{id}_Z\) iff \(c(\kappa_Z x,e)=\kappa_Z x\) for all \(x\) iff \(\kappa_Z\circ\tilde c=\kappa_Z\).
\end{proof}

The three descent conditions characterize exact quotient-level
composability. Associativity and identity are additional equations for
declared composites.

\paragraph{Nonclaims.} The theorem does not enumerate composable records and does not apply
to undeclared compositions.

\subsection*{Layer instantiations}\label{layer-inst:F28}

\begingroup\small
\begin{description}
\item[\emph{Universally composable cryptographic
protocols (theoretical computer science).}] \textit{Analogy.}
The carriers are implementations of protocol components;
the source quotients collapse implementations to their
UC ideal-functionality class; the interface readouts
record input-output behaviour visible to the
environment; the composition combines components by
message-passing; the composite quotient collapses the
composite to its ideal-functionality class.  The UC
composition theorem yields composite-outcome descent;
the residual records simulator output and descends
through the same component quotient.  Admissibility,
associativity, and identity correspond to syntactic
well-formedness of the protocol-composition operation
\citep{Canetti2001UC}.

\end{description}
\endgroup


\subsection{Presentation Invariance [\texorpdfstring{\Fref{F29}}{F29}]}\label{sec:status-records-coherence:presentation-invariance}

\paragraph{Reading the statement.}
Each predicate on the left is defined by the first condition displayed on
its right, as the symbol \(:\Longleftrightarrow\) marks: invariance under change of presentation within a
\(\pi\)-class, invariance under the declared compatible relation \(\sim\),
invariance of the relation \(R\) under the tuple quotient, and
covariance of \(T\) under the group.

\begin{theorem}[Presentation Invariance]\label{thm:F29}
Let \(P\) be a presentation space with quotient
\(\pi:P\to\Pi\), let \(\sim\) be a relation compatible
with \(\pi\), that is, \(p\sim p'\Rightarrow\pi(p)=\pi(p')\), let \(s:\Pi\to P\) be a section with
\(\pi(s(x))=x\), and let \(T:P\to V\) be a readout.  Let
\(\tau:\operatorname{Tuple}(P)\to\operatorname{Tuple}(\Pi)\) be the
tuple quotient on a declared carrier \(\operatorname{Tuple}(P)\) of presentation
tuples, for instance \(P^n\) with \(\tau=\pi^n\), let \(R:\operatorname{Tuple}(P)\to\mathrm{Prop}\) be a
relation, and let a group \(G\) act on \(P\) and \(V\).  Define
\begin{align*}
  \operatorname{PresentationInvariant}(\pi,T)
  &:\Longleftrightarrow
  \forall p,p'\in P,\ \pi(p)=\pi(p')\Rightarrow T(p)=T(p'),\\
  \operatorname{EquivalentPresentationInvariant}(\sim,T)
  &:\Longleftrightarrow
  \forall p,p'\in P,\ p\sim p'\Rightarrow T(p)=T(p'),\\
  \operatorname{PresentationRelationInvariant}(\tau,R)
  &:\Longleftrightarrow
  \forall x,y,\ \tau(x)=\tau(y)\Rightarrow (R(x)\leftrightarrow R(y)),\\
  \operatorname{GroupEquivariant}(T)
  &:\Longleftrightarrow
  \forall g\in G,\forall p\in P,\ T(g\cdot p)=g\cdot T(p).
\end{align*}
Then
\[
  \operatorname{PresentationInvariant}(\pi,T)
  \iff
  \exists\,\bar T:\Pi\to V\ \text{with}\ \bar T\circ\pi=T,
\]
\(\operatorname{PresentationInvariant}(\pi,T)\) implies
\(\operatorname{EquivalentPresentationInvariant}(\sim,T)\); when \(\sim\) is
an equivalence relation, the converse holds for every \(T\) into a
two-element \(V\) exactly when \(p\sim p'\iff\pi(p)=\pi(p')\); and
\(\operatorname{PresentationInvariant}(\pi,T)\) holds exactly when
\(T(p)=T(s(\pi(p)))\) for every \(p\in P\). More generally, if \(r\) is any
relation on \(P\) whose equivalence closure is the kernel of \(\pi\), that is,
\(\pi(p)=\pi(p')\) exactly when \(p\) and \(p'\) are joined by a finite chain
of \(r\)-steps taken in either direction, then
\(\operatorname{PresentationInvariant}(\pi,T)\) holds exactly when
\(T(p)=T(p')\) whenever \(r(p,p')\). For every \(\tau\), \(\operatorname{PresentationRelationInvariant}(\tau,R)\) holds if and only if
\(R=\bar R\circ\tau\) for some \(\bar R:\operatorname{Tuple}(\Pi)\to\mathrm{Prop}\).
\(\operatorname{GroupEquivariant}(T)\) is the displayed intertwining condition;
presentation invariance does not imply it.
\end{theorem}
\begin{proof}
Assume the quotient \(\pi\), compatible relation \(\sim\), section
\(s\), readout \(T\), tuple quotient \(\tau\), relation \(R\), and
group action of \(G\).  If \(T\) is presentation-invariant, it is
constant on \(\pi\)-fibers.  Conversely, if \(T\) is constant on
\(\pi\)-fibers, define
\[
  \bar T(x):=T(s(x)).
\]
For \(p\in P\), the presentations \(p\) and \(s(\pi(p))\) lie in the same
\(\pi\)-fiber, because \(\pi(s(\pi(p)))=\pi(p)\); since \(T\) is constant on
\(\pi\)-fibers, \(\bar T(\pi(p))=T(s(\pi(p)))=T(p)\), that is,
\(\bar T\circ\pi=T\). Since \(p\sim p'\) implies \(\pi(p)=\pi(p')\), a readout constant on
\(\pi\)-fibers is invariant under \(\sim\). If \(\sim\) is the kernel of
\(\pi\), the converse is immediate. If \(\sim\) is an equivalence relation
strictly smaller than the kernel, pick \(p\not\sim p'\) with
\(\pi(p)=\pi(p')\); the indicator of the \(\sim\)-class of \(p\) is
\(\sim\)-invariant but not presentation-invariant. If \(T\) is
presentation-invariant, then \(T(p)=T(s(\pi(p)))\) because \(p\) and
\(s(\pi(p))\) lie in one \(\pi\)-fiber; conversely, if \(T(p)=T(s(\pi(p)))\) for
every \(p\) and \(\pi(p)=\pi(p')\), then \(T(p)=T(s(\pi(p)))=T(s(\pi(p')))=T(p')\).

The relation statement holds for every \(\tau\): given invariance, put \(\bar R(z):\Leftrightarrow\exists x,\ \tau(x)=z\wedge R(x)\); then \(\bar R(\tau(y))\Leftrightarrow R(y)\). The converse is immediate.  The group-action statement is a separate
intertwining condition:
\[
  T(g\cdot p)=g\cdot T(p).
\]
For a generating relation \(r\): if \(T\) is presentation-invariant, then
\(r(p,p')\) gives \(\pi(p)=\pi(p')\) and so \(T(p)=T(p')\); conversely, if
\(T\) is constant along \(r\)-steps, it is constant along every finite chain of
steps taken in either direction, by induction on the length, hence on every
\(\pi\)-fiber.
For the non-implication, take \(P=\Pi=V=\{0,1\}\) and \(\pi=s=T=\mathrm{id}\),
and let \(C_2\) swap the points of \(P\) and act trivially on \(V\). Then \(T\) is
presentation-invariant, but \(T(g\cdot0)=1\ne0=g\cdot T(0)\) for the generator
\(g\).
\end{proof}

Readouts and relations satisfying the stated descent conditions are
independent of representatives of their respective presentation quotients.
The group-action clause characterizes equivariance; presentation independence
of a formed record requires descent of each relevant field.

\paragraph{Nonclaims.} The theorem does not enumerate presentations and does not collapse
the presentation/object distinction.

\subsection*{Layer instantiations}\label{layer-inst:F29}

\begingroup\small
\begin{description}
\item[\emph{Coordinate-free formulation of general
relativity (mathematical physics).}] \textit{Special case.}
The presentation space is the set of
coordinate-component representations of a geometric
tensor over a spacetime manifold; the quotient
collapses representations to the underlying
coordinate-free tensor field.  The diffeomorphism
group acts on representations by pushforward. The
readout \(T\) sending a coordinate representation to its tensor field is
presentation-invariant; it is also diffeomorphism-equivariant,
\(T(\phi\cdot p)=\phi_*T(p)\), which \Fref{F29} treats as a separate
condition.  General
covariance is the requirement that physically
meaningful readouts be presentation-invariant in this
sense \citep{MisnerThorneWheeler1973Gravitation}.

\item[\emph{Gr\"obner basis as canonical
polynomial-ideal presentation (computer algebra).}] \textit{Special case.}
The presentation space \(P\) is the set of finite subsets of
\(k[x_1,\dots,x_n]\), and \(\pi\) sends a set to the ideal it generates.
Under a fixed monomial ordering, the reduced Gr\"obner basis supplies a
canonical section: every ideal has a unique reduced Gr\"obner basis.
Membership of a fixed polynomial is a presentation-invariant readout.
Equality of generated ideals is an invariant relation on \(P^2\); for any
fixed ideal \(J\), so is the predicate \(I(A)\cap I(B)=J\), fitting the
relation clause with \(\tau=\pi^2\). Membership and equality can be checked
using Gr\"obner bases and normal forms, while an intersection can be computed
using an elimination Gr\"obner basis. Buchberger's algorithm constructs the
required bases \citep{Buchberger1965ThesisGroebner}.

\item[\emph{Jordan canonical form and matrix similarity
(linear algebra).}] \textit{Special case.}
For a finite-dimensional vector space \(V\) over an
algebraically closed field \(k\), the presentation
space is the set of matrix representatives of linear
endomorphisms; the quotient sends a matrix to its
similarity class.  Under a fixed eigenvalue and
block ordering, the Jordan canonical form supplies a
canonical section.  Trace, determinant, characteristic
polynomial, minimal polynomial, eigenvalue multiset
with algebraic and geometric multiplicities, rank, and
nullity are presentation-invariant readouts.  The
general linear group \(\mathrm{GL}_n(k)\) acts on
matrices by similarity; the listed invariants are
\(\mathrm{GL}_n(k)\)-invariant readouts, but they do not
separate similarity classes (nilpotent \(7\times 7\)
matrices with Jordan blocks of sizes \(3,3,1\) and
\(3,2,2\) agree on all of them), while the Jordan form,
read through the canonical section, is a complete
invariant \citep{HoffmanKunze1971LinearAlgebra}.

\item[\emph{De Bruijn indices for alpha-equivalence
(programming-language theory).}] \textit{Special case.}
The presentation space is the set of lambda terms with
named bound variables; the quotient sends a term to its
alpha-equivalence class.  The de Bruijn-index encoding
(each bound variable replaced by the distance to its
binder) is a complete invariant: two terms have the
same encoding iff they are alpha-equivalent.  Renaming
each bound variable by the depth of its binder, from a
reserved supply of names disjoint from the free
variables, gives a canonical named representative of
each class, which is the section \(s\).  The set of alpha-equivalence classes of one-step beta-reducts,
and the alpha-equivalence class of the beta-normal form when
one exists (with a marker otherwise), are presentation-invariant
readouts. Named-term outputs need not be presentation-invariant,
unless a canonical representative convention is imposed; the de Bruijn encoding
solves the variable-capture problem by construction
\citep{deBruijn1972LambdaNotation}.

\item[\emph{Lewin transformational music theory and the
chromatic group (music theory).}] \textit{Special case.}
The presentation space is the set of pitch-class
assignments to a musical motif; the chromatic group
\(G=\mathbb Z/12\mathbb Z\) acts by transposition, and
the dihedral group \(D_{12}\) of order 24 extends this with
inversion.  Some motif features, such as the pitch-class set, transform equivariantly,
\(T(g\cdot p)=g\cdot T(p)\); others are invariant.  Pitch-class invariants
such as the interval-class content (a canonical
\(G\)-invariant readout) descend through the orbit
quotient of the \(D_{12}\)-action and are
presentation-invariant musical features
\citep{Lewin1987GMIT}.
\end{description}
\endgroup

\subsection{Explanation as Compression with Descent [\texorpdfstring{\Fref{F30}}{F30}]}\label{sec:status-records-coherence:explanation}

\paragraph{Reading the statement.} The symbols are those of \Cref{thm:F30} below.
\(\operatorname{ExactExplanation}(b,e,t)\) says that the explanation
\(e\) descends through the base \(b\) and the target \(t\) descends
through \(e\): the base determines the explanation, and the explanation
determines the target. \(\operatorname{StrictCompression}(b,e)\) says
that \(e\) descends through \(b\) but \(b\) does not descend through
\(e\), so the explanation forgets some distinction of the base.

\begin{theorem}[Explanation as Compression with Descent]\label{thm:F30}
Let \(H\) be a history space with base readout \(b:H\to B\),
explanation readout \(e:H\to E\), target readout \(t:H\to T\),
continuation \(\kappa:H\to H'\), later explanation
\(e':H'\to E'\), and later target \(t':H'\to T'\).  Assume \(b\)
and \(e\) are surjective.  Define the obstruction sets
\begin{align*}
  \mathcal O_{\mathrm{src}}
  &:=
  \{(h,h')\in H^2:b(h)=b(h')\ \text{and}\ e(h)\ne e(h')\},\\
  \mathcal O_{\mathrm{tgt}}
  &:=
  \{(h,h')\in H^2:e(h)=e(h')\ \text{and}\ t(h)\ne t(h')\},\\
  \mathcal O_{\mathrm{cmp}}
  &:=
  \{(h,h')\in H^2:e(h)=e(h')\ \text{and}\ b(h)\ne b(h')\},
\end{align*}
and the predicates
\begin{align*}
  \operatorname{ExactExplanation}(b,e,t)
  &:\Longleftrightarrow
  \exists\,\alpha:B\to E,\ \exists\,\beta:E\to T,\
    \alpha\circ b=e\ \text{and}\ \beta\circ e=t,\\
  \operatorname{StrictCompression}(b,e)
  &:\Longleftrightarrow
  \exists\,\alpha:B\to E,\ \alpha\circ b=e
  \ \text{and}\ 
  \neg\exists\,\delta:E\to B,\ \delta\circ e=b,\\
  \operatorname{DynamicExplanation}(e,\kappa,e',t')
  &:\Longleftrightarrow
  \exists\,\gamma:E\to E',\ \gamma\circ e=e'\circ\kappa
  \ \text{and}\
  \exists\,\theta:E'\to T',\ \theta\circ e'=t'.
\end{align*}
Then:
\begin{enumerate}[label=(\roman*)]
\item \(\operatorname{ExactExplanation}(b,e,t)\) holds if and only if
  \(\mathcal O_{\mathrm{src}}=\emptyset\) and
  \(\mathcal O_{\mathrm{tgt}}=\emptyset\);
\item \(\operatorname{StrictCompression}(b,e)\) holds if and only if
  \(\mathcal O_{\mathrm{src}}=\emptyset\) and
  \(\mathcal O_{\mathrm{cmp}}\ne\emptyset\): the explanation is determined by
  the base and merges two histories that the base separates;
\item if \(\operatorname{DynamicExplanation}(e,\kappa,e',t')\), then the later
  target \(t'\circ\kappa\) descends through \(e\): the explanation now
  determines the target after the continuation.
\end{enumerate}
\end{theorem}
\begin{proof}
Assume the base readout \(b\), explanation readout \(e\), target
readout \(t\), and the surjectivity of \(b\) and \(e\).

(i) By the descent criterion of \Cref{sec:apparatus}, since \(b\) and
\(e\) are surjective, \(\alpha\) exists exactly when
\(\mathcal O_{\mathrm{src}}=\emptyset\), and \(\beta\) exactly when
\(\mathcal O_{\mathrm{tgt}}=\emptyset\).

(ii) By (i), \(\alpha\) exists exactly when
\(\mathcal O_{\mathrm{src}}=\emptyset\), and by the descent criterion for
the surjective \(e\), a map \(\delta:E\to B\) with \(\delta\circ e=b\)
exists exactly when \(\mathcal O_{\mathrm{cmp}}=\emptyset\).  So strict compression is
\(\mathcal O_{\mathrm{src}}=\emptyset\) together with
\(\mathcal O_{\mathrm{cmp}}\ne\emptyset\).

(iii) Given \(\gamma\) and \(\theta\),
\(t'\circ\kappa=\theta\circ e'\circ\kappa=\theta\circ\gamma\circ e\), so
\(t'\circ\kappa\) factors through \(e\).
\end{proof}

\paragraph{Nonclaims.} The theorem does not rank explanations and does not identify
explanation with the fact explained.  The Kolmogorov-complexity and
minimum-description-length traditions develop the
shortest-program / shortest-code reading of explanation in detail
\citep{LiVitanyi2008Kolmogorov}.

\subsection*{Layer instantiations}\label{layer-inst:F30}

\begingroup\small
\begin{description}
\item[\emph{Newtonian mechanics explaining
Kepler's laws (mathematical physics).}] \textit{Analogy.}
The history space is the set of solar-system
trajectories; the base records observed positions and
velocities.  The explanation is Newton's gravitational
dynamics: masses of the bodies, the inverse-square
gravitational law, and initial conditions.  The target
is the set of Kepler-type empirical regularities
(elliptical orbits, equal areas in equal times,
\(T^2\propto a^3\)) and the ephemeris of future
planetary positions.  The descended map identifies
masses and initial conditions from kinematics, and
the target map solves the two-body problem to derive
Kepler's laws; because the dynamics is deterministic,
integrating the equations of motion recovers the trajectory from the
explanation, so this compression is exact but not strict in the
\Fref{F30} sense; its gain is that finitely many parameters and one
law replace the continuous trajectory \citep{Newton1687Principia}.

\item[\emph{Principal component analysis as
dimensionality-reducing explanation (statistics).}] \textit{Special case.}
The history space is observed records
\(x\in\mathbb R^p\); the base is the full record.
The explanation is the top-\(k\) principal-component
representation, computed as eigenvectors of the
empirical covariance matrix.  The target is the
low-rank reconstruction or rank-\(k\) orthogonal
projection visualisation, both exactly determined by
the principal-component representation through a
linear map.  Taking the history space to be all of
\(\mathbb R^p\), the compression is strict whenever \(k<p\), because
the rank-\(k\) projection is not injective; on a finite data set the
projection can be injective, and non-zero discarded variance then
measures reconstruction error rather than strictness.  The
explained-variance ratio measures the compression fidelity \citep{Pearson1901PCA}.

\end{description}
\endgroup


\subsection{Counterfactual Legitimacy [\texorpdfstring{\Fref{F31}}{F31}]}\label{sec:status-records-coherence:counterfactual}

\paragraph{Reading the statement.} The symbols are those of \Cref{thm:F31} below.
\(\operatorname{CFL}\) says that admissibility is constant on
\((q,\alpha)\)-classes and that target truth is constant on the
admissible pairs of each class, so that a counterfactual claim depends
only on the declared context and intervention classes. The target
quotient \(q'\) is recorded with the data, but
neither obstruction set uses it.

\begin{theorem}[Counterfactual Legitimacy]\label{thm:F31}
Let \(q:H\to Q\) be a context quotient, let \(\alpha:A\to\bar A\) be an intervention abstraction, let \(\operatorname{Adm}:H\times A\to\mathrm{Prop}\) be admissibility, let
\(\Gamma:H\times A\to H'\) be a counterfactual route, let
\(q':H'\to R\) be a target quotient, and let
\(\tau:H'\to\mathrm{Bool}\) be target truth.  Define
\[
  \mathcal O_{\mathrm{adm}}
  :=
  \{((h,a),(h',a')) :
    q(h)=q(h'),\ \alpha(a)=\alpha(a'),\
    \operatorname{Adm}(h,a)\not\leftrightarrow
    \operatorname{Adm}(h',a')\},
\]
and
\[
  \mathcal O_{\mathrm{cf}}
  :=
  \left\{
  \begin{array}{l}
    ((h,a),(h',a')) :
    q(h)=q(h'),\ \alpha(a)=\alpha(a'),\\
    \operatorname{Adm}(h,a),\ \operatorname{Adm}(h',a'),\\
    \tau(\Gamma(h,a))\ne\tau(\Gamma(h',a'))
  \end{array}
  \right\}.
\]
Define the counterfactual-legitimacy predicate by
\(\operatorname{CFL}(q,\alpha,\operatorname{Adm},\Gamma,q',\tau):\Longleftrightarrow\)
for all \(h,h'\in H\) and \(a,a'\in A\) with \(q(h)=q(h')\) and
\(\alpha(a)=\alpha(a')\), \(\operatorname{Adm}(h,a)\leftrightarrow\operatorname{Adm}(h',a')\),
and if moreover \(\operatorname{Adm}(h,a)\) and \(\operatorname{Adm}(h',a')\), then
\(\tau(\Gamma(h,a))=\tau(\Gamma(h',a'))\); the target quotient \(q'\) is not used
in this definition.  Then
\begin{align*}
  \operatorname{CFL}(q,\alpha,\operatorname{Adm},\Gamma,q',\tau)
  &\iff
  \mathcal O_{\mathrm{adm}}=\emptyset
  \wedge
  \mathcal O_{\mathrm{cf}}=\emptyset\\
  &\iff
  \begin{gathered}
    \exists\,\overline{\operatorname{Adm}}:Q\times\bar A\to\mathrm{Prop},\
    \exists\,\bar\tau:Q\times\bar A\to\mathrm{Bool}
    \ \text{such that}\\
    \overline{\operatorname{Adm}}(q(h),\alpha(a))
      \leftrightarrow\operatorname{Adm}(h,a)
    \ \text{for all }(h,a)\ \text{and}\\
    \bar\tau(q(h),\alpha(a))=\tau(\Gamma(h,a))
    \ \text{for all admissible }(h,a).
  \end{gathered}
\end{align*}
\end{theorem}
\begin{proof}
Assume the context quotient \(q\), intervention abstraction
\(\alpha\), admissibility predicate \(\operatorname{Adm}\), route
\(\Gamma\), target quotient \(q'\), and target truth \(\tau\).  If a
counterfactual is legitimate, admissibility is constant on
\((q,\alpha)\)-classes, so
\(\mathcal O_{\mathrm{adm}}=\emptyset\).  Its target truth is also
constant on admissible \((q,\alpha)\)-classes, so
\(\mathcal O_{\mathrm{cf}}=\emptyset\).

Conversely, assume both obstruction sets are empty.  Define
\[
  \overline{\operatorname{Adm}}(q(h),\alpha(a)):=\operatorname{Adm}(h,a)
\]
and, for admissible pairs,
\[
  \bar\tau(q(h),\alpha(a)):=\tau(\Gamma(h,a)).
\]
The emptiness of \(\mathcal O_{\mathrm{adm}}\) makes
\(\overline{\operatorname{Adm}}\) well-defined, it is set false off the image of \(q\times\alpha\), \(\bar\tau\) is set to false at classes with no admissible representative, and the emptiness of
\(\mathcal O_{\mathrm{cf}}\) makes \(\bar\tau\) independent of the
admissible representative.  Conversely again, given
\(\overline{\operatorname{Adm}}\) and \(\bar\tau\), two pairs with the
same classes have the same admissibility and, when admissible, the
same target truth, so both obstruction sets are empty.  Thus a pair in either obstruction set is exactly a witness that
admissibility or admissible target truth fails to descend to
\(Q\times\bar A\).
\end{proof}

\paragraph{Nonclaims.} The theorem does not determine the truth value of counterfactuals and
does not apply to arbitrary undeclared possible worlds.  The
possible-worlds semantics in the background is the classical reference
\citep{Lewis1973Counterfactuals}.

\subsection*{Layer instantiations}\label{layer-inst:F31}

\begingroup\small
\begin{description}
\item[\emph{{MIL-STD-810} environmental-qualification
bench testing (engineering qualification).}] \textit{Analogy.}
US Department of Defense MIL-STD-810 prescribes a
catalogue of environmental qualification test methods
(thermal cycling, vibration, shock, humidity, salt fog,
etc.).  The context quotient collapses specimens to
their design baseline (model number, build
configuration); the intervention abstraction collapses
test profiles to the standardised method class;
admissibility requires the profile to lie within the
standard envelope; the counterfactual route is the
laboratory test execution under the prescribed
profile; the target truth records pass/fail according
to the acceptance criterion.  The legitimacy of a
qualification-test counterfactual --- ``this specimen
would pass under method X'' --- is the \Fref{F31} statement:
admissibility and pass/fail descend through
(design-baseline, test-method) classes
\citep{MILSTD810H2019}.

\end{description}
\endgroup


\subsection*{Status, Records, Coherence at a glance}\label{sec:status-records-coherence:summary}\addcontentsline{toc}{subsection}{Status, Records, Coherence at a glance}
\begin{center}
  \centering
  \footnotesize
  \setlength{\tabcolsep}{4pt}
  \medskip\noindent\begin{minipage}{\textwidth}\centering
  \textbf{Status}\par\smallskip
  \begin{tabularx}{\textwidth}{@{}
      >{\raggedright\arraybackslash}p{0.24\textwidth}
      >{\raggedright\arraybackslash}X
      >{\raggedright\arraybackslash}p{0.10\textwidth}@{}}
    \toprule
    Name & One-line claim & Substrate cite \\
    \midrule
    Objectivity as Common Quotient (\Fref{F17}) &
    For surjective access and public quotients, when the public quotient is the common quotient of the accesses (it factors through each, and every map factoring through each factors through it), a readout factors through every access exactly when the public quotient determines it, and exactly when no access identifies two histories with different readout values.  &
    \textemdash \\
    Lawhood-Descent Theorem (\Fref{F18}) &
    For a surjective access quotient, a candidate law is lawful exactly
    when it factors through it.  &
    \textemdash \\
    Interface-Mediation Theorem (\Fref{F22}) &
    For surjective source quotient and boundary readout, an update is mediated iff its target class descends, uniquely, through (source
    class, boundary value); failure splits into boundary leakage and internal
    insufficiency.  &
    \textemdash \\
    Probability as Stable Unresolved Fiber (\Fref{F23}) &
    For a surjective quotient, a readout either descends or is unresolved, never both; stable
    probability is unresolvedness together with the declared ledger
    conditions Stable, which do not involve the weights and hold trivially for empty supports and a one-point target set; on a finite carrier, with the weight assignment given on fiber events by real fiber
    weights, positive and summing to one on each fiber, descent holds iff every fiber chance is a point mass (value 1 at one readout value).  &
    \textemdash \\
    \bottomrule
  \end{tabularx}
  \end{minipage}\par

  \medskip\noindent\begin{minipage}{\textwidth}\centering
  \textbf{Records}\par\smallskip
  \begin{tabularx}{\textwidth}{@{}
      >{\raggedright\arraybackslash}p{0.24\textwidth}
      >{\raggedright\arraybackslash}X
      >{\raggedright\arraybackslash}p{0.10\textwidth}@{}}
    \toprule
    Name & One-line claim & Substrate cite \\
    \midrule
    Object-Persistence Theorem (\Fref{F19}) &
    For a surjective source access quotient \(q_i\), an object persists exactly when its split obstruction is empty; exactly one of four comparison statuses holds; the pairing \(h\mapsto(q_i(h),q_j(\gamma h))\) retains \(q_i\) and carries \(q_j\circ\gamma\); viability, relations and holonomy hold exactly when their violation sets are empty.  &
    \textemdash \\
    Fact--Record Theorem (\Fref{F20}) &
    For surjective record signatures \(s\) and \(\rho_s\), an event is recorded, a record is stable, and a fact value is stable
    exactly when the corresponding obstruction is empty; the record
    comparison has exactly one of four statuses.  &
    \textemdash \\
    Reflexive-Nonclosure Theorem (\Fref{F21}) &
    Under rules 2--4 of a same-layer audit policy, each of which is needed, a complete self-soundness claim is never accepted, and, when the meta layer accepts no self-sound meta-claim, a meta-layer certificate comes from an accepted
    meta-claim that is not itself self-sound.  &
    \textemdash \\
    Causality Theorem (\Fref{F25}) &
    For surjective context and intervention quotients, intervention
    transport descends to (context class, intervention
    class) exactly when its obstruction is empty, and then uniquely;
    under a descended transport every stable causal effect has a witness
    at one raw context with two raw interventions.  &
    \textemdash \\
    \bottomrule
  \end{tabularx}
  \end{minipage}\par

  \medskip\noindent\begin{minipage}{\textwidth}\centering
  \textbf{Coherence}\par\smallskip
  \begin{tabularx}{\textwidth}{@{}
      >{\raggedright\arraybackslash}p{0.24\textwidth}
      >{\raggedright\arraybackslash}X
      >{\raggedright\arraybackslash}p{0.10\textwidth}@{}}
    \toprule
    Name & One-line claim & Substrate cite \\
    \midrule
    Moduli--Contingency Theorem (\Fref{F26}) &
    When the moduli map identifies exactly the equivalent lawful closures and every modulus is attained by a lawful closure, a readout agrees
    on lawful closures with a map on the moduli set
    exactly when it is presentation-independent on lawful closures; it is then either
    necessary or contingent, never both.  &
    \textemdash \\
    Conservation Theorem (\Fref{F27}) &
    Conserved quantities are exactly readouts that factor through
    route-orbit quotients.  &
    \textemdash \\
    Composability Theorem (\Fref{F28}) &
    For a surjective part/interface quotient \(\kappa\), the composite outcome, cross residual and admissibility each descend through \(\kappa\), uniquely, exactly when their split-pair sets are empty; exact composability is emptiness of all three.  &
    \textemdash \\
    Presentation-Invariance Theorem (\Fref{F29}) &
    For a quotient with a section, a readout constant on presentation classes is exactly one that factors through the quotient, agrees with its value at the section representative, or is constant along any relation generating the kernel; relation invariance under any tuple quotient is descent; presentation invariance does not imply equivariance.  &
    \textemdash \\
    Explanation-as-Compression Theorem (\Fref{F30}) &
    For surjective base and explanation readouts, an explanation is exact
    exactly when its source and target obstructions are empty, and is a
    strict compression exactly when it is determined by the base and merges
    two histories the base separates; a dynamic explanation makes the later
    target descend through the explanation.  &
    \textemdash \\
    Counterfactual-Legitimacy Theorem (\Fref{F31}) &
    A counterfactual is legitimate exactly when both obstruction sets are empty, exactly when admissibility, and target truth on admissible pairs, descend to (context class, intervention class).  &
    \textemdash \\
    \bottomrule
  \end{tabularx}
  \end{minipage}\par

  \captionof{table}{Status, Records, Coherence at a glance: the fourteen laws of \Cref{sec:status-records-coherence}, grouped into reading groups.}
  \label{tab:axis-status-records-coherence-summary}
\end{center}

\FloatBarrier
